% =============================================================
%  MÉMOIRE DE PROJET DE FIN D'ÉTUDES
%  EDURA -- Plateforme SaaS de Gestion Scolaire Intelligente
%  + Volet Entrepreneurial (dispositif Diplôme-Startup)
%  Amina Kaouter Ghezal -- USTO-MB -- 2025/2026
%
%  Compilation : XeLaTeX ou pdfLaTeX (Overleaf : XeLaTeX recommandé)
%  Bibliographie : biber  (dans Overleaf : Settings > Compiler > biber)
% =============================================================

\documentclass[12pt,a4paper,oneside]{report}
\usepackage[utf8]{inputenc}

% ---- Encodage & langue ----
\usepackage[utf8]{inputenc}
\usepackage[T1]{fontenc}
\usepackage[french]{babel}
\usepackage{lmodern}

% ---- Mise en page ----
\usepackage[top=3cm,bottom=2.5cm,left=3cm,right=2.5cm]{geometry}
\usepackage{setspace}
\onehalfspacing
\usepackage{parskip}
\setlength{\parskip}{6pt}

% ---- En-têtes & pieds de page ----
\usepackage{fancyhdr}
\setlength{\headheight}{15pt}
\pagestyle{fancy}
\fancyhf{}
\fancyhead[L]{\small\leftmark}
\fancyhead[R]{\small EDURA -- USTO-MB 2025/2026}
\fancyfoot[C]{\thepage}
\renewcommand{\headrulewidth}{0.4pt}

% ---- Titres de chapitres ----
\usepackage{titlesec}
\titleformat{\chapter}[display]
  {\normalfont\huge\bfseries}
  {\chaptertitlename\ \thechapter}{18pt}{\Huge}
\titlespacing*{\chapter}{0pt}{30pt}{20pt}

% ---- Graphiques & figures ----
\usepackage{graphicx}
\usepackage{float}
\usepackage{caption}
\usepackage{subcaption}
\usepackage[dvipsnames]{xcolor}
\graphicspath{{./figures/}}

% ---- Tableaux ----
\usepackage{booktabs}
\usepackage{array}
\usepackage{longtable}
\usepackage{multirow}
\usepackage{tabularx}
\usepackage{colortbl}

% ---- Code source ----
\usepackage{listings}
\lstset{
  basicstyle=\ttfamily\footnotesize,
  keywordstyle=\color{blue}\bfseries,
  commentstyle=\color{gray},
  stringstyle=\color{red!70!black},
  breaklines=true,
  frame=single,
  numbers=left,
  numberstyle=\tiny\color{gray},
  captionpos=b,
  tabsize=2,
  showstringspaces=false
}

% ---- Maths ----
\usepackage{amsmath}
\usepackage{amssymb}

% ---- Boîtes ----
\usepackage{mdframed}
\usepackage{csquotes}

\usepackage[backend=biber,style=ieee,sorting=none,doi=true,url=false]{biblatex}
\newbibmacro*{volume+number}{%
  \printfield{volume}%
  \setunit*{\addnbspace}%
  \printfield{number}%
}
\addbibresource{references.bib}

% ---- Hyperliens ----
\usepackage[colorlinks=true,linkcolor=blue!70!black,citecolor=green!50!black,urlcolor=blue]{hyperref}

% ---- Table des matières ----
\usepackage{tocloft}
\renewcommand{\cftchapfont}{\bfseries}
\setlength{\cftbeforechapskip}{6pt}

% ---- Acronymes ----
\usepackage{acronym}

% ---- Commandes personnalisées ----
\newcommand{\edura}{\textsc{Edura}}
\newcommand{\todo}[1]{\textcolor{red}{\textbf{[TODO: #1]}}}

% Boîte d'info pour les explications
\newmdenv[backgroundcolor=gray!10,linecolor=gray!40,
          innerleftmargin=10pt,innerrightmargin=10pt,
          innertopmargin=6pt,innerbottommargin=6pt]{infobox}

% Cadre placeholder pour les captures d'écran à insérer
\newcommand{\screenplaceholder}[2]{%
  \begin{figure}[H]
    \centering
    \setlength{\fboxsep}{12pt}
    \setlength{\fboxrule}{0.5pt}
    \fbox{\parbox{0.85\textwidth}{\centering\vspace{2cm}
    \textit{[CAPTURE D'ÉCRAN À INSÉRER :}\\[6pt]
    \textbf{#1}]
    \vspace{2cm}}}
    \caption{#1}
    \label{fig:#2}
  \end{figure}
}

% =============================================================
\begin{document}
% =============================================================

% ---- PAGE DE GARDE ----
\clearpage

% ---- DÉDICACES ----
\clearpage
\thispagestyle{empty}
\vspace*{4cm}
\begin{flushright}
\itshape
À mes parents, mes soeurs, pour leur soutien indéfectible\\
tout au long de ce parcours.\\[12pt]
À tous ceux qui croient que l'éducation\\
peut et doit mieux faire pour chaque enfant.
\end{flushright}
\clearpage

% ---- REMERCIEMENTS ----
\chapter*{Remerciements}
\addcontentsline{toc}{chapter}{Remerciements}
\thispagestyle{plain}

Je tiens d'abord à exprimer ma sincère gratitude à Madame \textsc{Dekhici} Latifa, mon encadrante, pour ses conseils éclairés, sa disponibilité et la liberté qu'elle m'a accordée dans l'exploration de sujets à la croisée de l'informatique et des sciences de l'éducation. Ses remarques ont considérablement affiné la rigueur de ce travail.

Mes remerciements vont également aux membres du jury, qui ont accepté d'évaluer ce mémoire avec le regard critique que mérite un tel travail.

Je remercie l'ensemble du corps enseignant du département d'Informatique de l'USTO-MB pour la formation solide dispensée au cours de ces cinq années. Les fondements théoriques acquis dans cet établissement ont directement nourri les choix d'architecture décrits dans ce document.

Je remercie l'incubateur universitaire de l'USTO-MB et les responsables du dispositif Diplôme-Startup, qui ont enrichi ce projet d'une dimension entrepreneuriale en plus de sa dimension technique. Le guide méthodologique en six axes qui structure le chapitre entrepreneurial a permis de transformer une démonstration technique en un véritable plan de création d'entreprise.

Enfin, je remercie les directeurs et secrétaires des écoles privées d'Oran qui ont bien voulu me consacrer du temps lors de l'étude qualitative qui a précédé le développement de la plateforme. Leurs témoignages ont été déterminants pour comprendre les véritables contraintes du terrain.

% ---- RÉSUMÉ ----
\chapter*{Résumé}
\addcontentsline{toc}{chapter}{Résumé}
\thispagestyle{plain}

\noindent\textbf{Titre :} Conception et réalisation d'une plateforme SaaS de gestion pour établissements scolaires : analyse prédictive des risques d'échec, profilage psychopédagogique (Gardner et MBTI) et orientation automatisée des élèves.

\vspace{8pt}

Le secteur de l'enseignement privé algérien, qui compte aujourd'hui plus de 800 établissements actifs, souffre d'une double lacune : une gestion administrative fragmentée, essentiellement manuelle, et l'absence totale d'outils permettant de valoriser scientifiquement le potentiel cognitif de chaque élève. Ce mémoire présente \edura{}, une plateforme SaaS multi-tenant conçue pour répondre simultanément à ces deux problèmes.

Sur le plan administratif, la plateforme centralise treize modules fonctionnels couvrant la gestion des élèves, des notes, des présences, de la finance, de la conformité réglementaire et de la communication interne. Sur le plan pédagogique, elle intègre une couche d'intelligence artificielle hybride articulée autour de deux paradigmes complémentaires : l'IA symbolique, sous forme de systèmes experts à règles pondérées issus de la littérature en sciences de l'éducation, et l'IA générative via l'API Anthropic Claude Haiku pour la production de résumés en langage naturel.

Le module phare d'\edura{} est le \textbf{Rapport Scientifique d'Orientation}, qui opérationnalise la théorie des intelligences multiples de Howard Gardner \autocite{gardner1983frames} et l'indicateur de personnalité Myers-Briggs (MBTI) \autocite{myers1980gifts} à partir des données scolaires standards, pour produire pour chaque élève une prédiction de filière BAC, un profil psychopédagogique complet et des recommandations personnalisées de métiers et d'universités algériennes. Une application mobile compagnon -- \textit{EDURA Test App}, développée en React Native -- permet aux élèves de passer les tests d'intelligences multiples, de MBTI et d'orientation directement sur une tablette, les résultats étant automatiquement synchronisés avec la plateforme principale.

La plateforme est développée avec Next.js, TypeScript, PostgreSQL et Supabase, et est déployée en production sur infrastructure Vercel. Elle a été livrée sous forme de MVP fonctionnel en six mois de développement solo, parallèlement à la formation académique. Le présent mémoire intègre également un \textbf{volet entrepreneurial complet} structuré en six axes (présentation, innovation, marché, production, finances, prototype) conformément au dispositif Diplôme-Startup.

\vspace{10pt}
\noindent\textbf{Mots-clés :} SaaS, multi-tenant, gestion scolaire, intelligence artificielle, décrochage scolaire, intelligences multiples, MBTI, orientation pédagogique, Next.js, PostgreSQL, React Native, Algérie, entrepreneuriat étudiant.

\vspace{20pt}
\noindent\rule{\linewidth}{0.4pt}

\vspace{10pt}
\noindent\textbf{Abstract}

\vspace{6pt}
Algeria's private education sector, with over 800 active institutions, suffers from two structural weaknesses: fragmented, mostly paper-based administrative management and a complete absence of tools capable of scientifically assessing each student's cognitive potential. This thesis presents \edura{}, a multi-tenant SaaS platform designed to address both issues simultaneously.

On the administrative side, the platform consolidates thirteen functional modules covering student management, grades, attendance, finance, regulatory compliance, and internal communication. On the pedagogical side, it integrates a hybrid AI layer combining rule-based expert systems grounded in educational research literature with generative AI via the Anthropic Claude Haiku API for natural-language summary generation.

\edura{}'s flagship feature is the Individual Scientific Orientation Report, which operationalises Howard Gardner's theory of multiple intelligences and the Myers-Briggs Type Indicator from standard academic data, producing for each student a baccalaureate track prediction, a complete psychopedagogical profile, and personalised career and university recommendations. A companion React Native tablet app -- \textit{EDURA Test App} -- enables students to take the intelligence, MBTI and career-orientation tests directly on a tablet, with results automatically synchronised back to the main platform.

The platform is built with Next.js, TypeScript, PostgreSQL, and Supabase, deployed in production on Vercel infrastructure. It was delivered as a functional MVP within six months of solo development, concurrent with academic studies. This thesis also includes a complete \textbf{entrepreneurial component} structured in six axes (presentation, innovation, market, production, finance, prototype) compliant with the national Diplôme-Startup framework.

\vspace{10pt}
\noindent\textbf{Keywords:} SaaS, multi-tenancy, school management, artificial intelligence, dropout risk, multiple intelligences, MBTI, academic orientation, Next.js, PostgreSQL, React Native, Algeria, student entrepreneurship.

\clearpage
% ---- TABLE DES MATIÈRES ----
\tableofcontents

% ---- LISTE DES FIGURES ----
\listoffigures
\addcontentsline{toc}{chapter}{Liste des figures}

% ---- LISTE DES TABLEAUX ----
\listoftables
\addcontentsline{toc}{chapter}{Liste des tableaux}

% ---- LISTE DES ACRONYMES ----
\chapter*{Liste des acronymes}
\addcontentsline{toc}{chapter}{Liste des acronymes}
\begin{acronym}[USTO-MB]
  \acro{SaaS}{Software as a Service}
  \acro{ERP}{Enterprise Resource Planning}
  \acro{IA}{Intelligence Artificielle}
  \acro{GOFAI}{Good Old-Fashioned Artificial Intelligence}
  \acro{LLM}{Large Language Model}
  \acro{ML}{Machine Learning}
  \acro{MBTI}{Myers-Briggs Type Indicator}
  \acro{RBAC}{Role-Based Access Control}
  \acro{JWT}{JSON Web Token}
  \acro{RLS}{Row-Level Security}
  \acro{ORM}{Object-Relational Mapping}
  \acro{CI/CD}{Continuous Integration / Continuous Deployment}
  \acro{SSR}{Server-Side Rendering}
  \acro{API}{Application Programming Interface}
  \acro{RGPD}{Règlement Général sur la Protection des Données}
  \acro{MVP}{Minimum Viable Product}
  \acro{BAC}{Baccalauréat}
  \acro{USTO-MB}{Université des Sciences et de la Technologie d'Oran -- Mohamed Boudiaf}
  \acro{DZD}{Dinar Algérien}
  \acro{BMC}{Business Model Canvas}
  \acro{BFR}{Besoin en Fonds de Roulement}
  \acro{CA}{Chiffre d'Affaires}
  \acro{TAM}{Total Addressable Market}
  \acro{SAM}{Serviceable Addressable Market}
  \acro{SOM}{Serviceable Obtainable Market}
\end{acronym}

% =============================================================
%  INTRODUCTION GÉNÉRALE
% =============================================================
\chapter*{Introduction générale}
\addcontentsline{toc}{chapter}{Introduction générale}
\markboth{Introduction générale}{Introduction générale}

L'école privée algérienne occupe aujourd'hui une place singulière dans le paysage éducatif national. Depuis que le décret exécutif n°~04-90 du 24 mars 2004 a encadré juridiquement leur création \autocite{decret2004}, ces établissements ont connu une croissance remarquable, passant de quelques dizaines à plus de 800 structures actives réparties sur l'ensemble du territoire. Pourtant, malgré cette expansion, les outils de gestion dont disposent leurs directeurs demeurent d'une pauvreté frappante : des tableurs Excel disjoints, des cahiers de présence en papier, des messages WhatsApp pour communiquer avec les parents. Dans un secteur où la moindre défaillance administrative peut entraîner la fermeture d'un établissement -- vingt écoles ont ainsi été sanctionnées depuis 2023 pour non-conformité réglementaire --, cette fragilité structurelle est difficile à justifier.

Ce constat, établi au fil d'entretiens menés auprès de directeurs et de secrétaires d'écoles privées oranaises, constitue le point de départ de ce mémoire. Il soulève une question à la fois technique et sociale : est-il possible de concevoir, en parallèle d'une formation académique et dans un contexte de ressources limitées, une plateforme numérique réellement adaptée aux contraintes des établissements privés algériens ? La réponse que propose ce travail est affirmative, mais elle impose des choix clairs sur ce qu'il faut construire et ce qu'il faut, pour l'instant, différer.

\edura{} est le nom de cette plateforme. Il s'agit d'une application web de type SaaS à architecture multi-tenant, développée intégralement en TypeScript avec le framework Next.js, et déployée en production à l'adresse \url{https://edura-aminaghezal.vercel.app}. Elle s'articule autour de deux axes complémentaires : la rationalisation de la gestion administrative scolaire d'une part, et la valorisation scientifique du potentiel individuel de l'élève de l'autre. Ce second axe, à travers le \textit{Rapport Scientifique d'Orientation}, constitue la contribution la plus originale du projet. Le rapport mobilise simultanément la théorie des intelligences multiples de Howard Gardner et l'indicateur de personnalité Myers-Briggs (\ac{MBTI}), couplée à une application mobile compagnon en React Native qui permet aux élèves de passer les tests psychométriques directement sur une tablette.

\vspace{6pt}
\noindent\textbf{Organisation du mémoire.} Ce document est structuré en six chapitres et un volet entrepreneurial.

Le premier chapitre pose le cadre du marché et dresse un état des généralités nécessaires à la compréhension du projet. Le deuxième recense les travaux connexes et les solutions existantes, tant en matière de gestion scolaire que d'orientation pédagogique assistée. Le troisième chapitre décrit la démarche de conception : architecture logicielle, modèle de données, choix technologiques. Le quatrième chapitre détaille la réalisation effective des modules, incluant le module phare du Rapport Scientifique enrichi (Gardner + MBTI + photo élève) et l'application compagnon \textit{EDURA Test App}. Le cinquième présente les tests conduits et les résultats obtenus.

Le sixième chapitre, ajouté en réponse au dispositif Diplôme-Startup et au guide méthodologique fourni par l'incubateur universitaire, présente le \textbf{volet entrepreneurial} structuré en six axes : présentation du projet, aspects innovants, analyse stratégique du marché, plan de production et organisation, plan financier, et prototype expérimental.

% =============================================================
%  CHAPITRE 1 : ÉTUDE DU MARCHÉ ET GÉNÉRALITÉS
% =============================================================
\chapter{Étude du marché et généralités}

\section{L'enseignement privé en Algérie : un secteur en mutation}

\subsection{Panorama chiffré}

L'autorisation des établissements d'enseignement privés en Algérie remonte officiellement à 2004, avec la promulgation du décret exécutif n°~04-90 \autocite{decret2004}. En moins de vingt ans, le nombre d'écoles a progressé de façon régulière jusqu'à atteindre une estimation de 800 structures actives sur le territoire national, concentrées principalement dans les wilayas d'Alger, d'Oran, de Constantine, de Sétif et d'Annaba. Ces chiffres, bien qu'il faille les manier avec prudence en l'absence de données officielles consolidées, témoignent d'une demande sociale réelle : des familles algériennes cherchent des cadres scolaires plus encadrés, plus flexibles ou simplement différents de l'offre publique.

Sur le plan économique, les frais de scolarité pratiqués varient considérablement d'un établissement à l'autre, allant de 200 000 à 600 000 dinars algériens (\ac{DZD}) par élève et par an. Un établissement de taille moyenne, soit entre 150 et 400 élèves, génère donc un chiffre d'affaires annuel compris entre 30 et 240 millions \ac{DZD}. C'est suffisant pour justifier un investissement dans des outils numériques sérieux ; c'est en même temps assez modeste pour rendre prohibitifs les logiciels internationaux dont les licences sont libellées en dollars ou en euros.

\subsection{Les fragilités structurelles identifiées}

Une étude qualitative conduite auprès d'une dizaine de directeurs et secrétaires d'écoles privées oranaises avant le démarrage du développement a permis d'identifier six problèmes récurrents. Ils méritent d'être décrits avec un minimum de détail, car ils ont directement orienté les choix fonctionnels d'\edura{}.

Le plus visible est la fragmentation des données opérationnelles. Aucun de ces établissements ne disposait d'un système centralisé : les notes étaient dans des fichiers Excel (souvent un fichier par classe et par trimestre, tenus par chaque enseignant selon ses propres conventions), les présences dans des cahiers papier, les paiements dans un tableau différent. Il n'existait nulle part de vue d'ensemble. La conséquence pratique en est que préparer un conseil de classe demande plusieurs heures de ressaisie manuelle, avec les erreurs que cela implique.

Plus grave encore, aucun de ces établissements ne pratiquait la détection précoce du décrochage scolaire. La littérature académique est pourtant formelle : c'est l'absentéisme et la chute brutale des résultats, détectés suffisamment tôt, qui permettent une intervention efficace \autocite{balfanz2007preventing}. Dans les écoles observées, les élèves en difficulté n'étaient identifiés qu'à la remise des bulletins trimestriels, soit souvent bien trop tard.

La troisième fragilité concerne la conformité réglementaire. Le cahier des charges des établissements privés algériens, mis à jour en 2026, impose des obligations précises sur huit catégories : qualifications des enseignants, respect du programme officiel, conditions d'infrastructure, sécurité des locaux, et autres. La non-conformité peut entraîner une fermeture administrative, comme en témoignent les vingt établissements sanctionnés depuis 2023. Or, aucun des directeurs interrogés ne disposait d'un outil pour suivre ces obligations en temps réel.

\subsection{Le problème de fond : l'absence de valorisation du potentiel individuel}

Il existe un problème qui transcende la gestion quotidienne et qui constitue, d'un point de vue social, la lacune la plus préoccupante. En Algérie, comme dans de nombreux pays du Maghreb, l'orientation scolaire en fin de cycle secondaire repose quasi exclusivement sur les moyennes générales. L'élève qui obtient la meilleure note est orienté vers la filière « Sciences expérimentales » ou « Mathématiques », indépendamment de ses aptitudes réelles, de son style d'apprentissage, de ses intérêts ou de ses intelligences dominantes.

Cette pratique est en contradiction directe avec ce que la recherche en psychologie cognitive indique depuis plus de quarante ans. Howard Gardner, psychologue à l'Université Harvard, a montré dès 1983 que l'intelligence humaine n'est pas une capacité uniforme mesurable par un seul chiffre, mais un ensemble de huit aptitudes distinctes qui se combinent différemment chez chaque individu \autocite{gardner1983frames}. De même, l'indicateur Myers-Briggs (MBTI), opérationnalisation moderne des types psychologiques de Carl Jung \autocite{myers1980gifts}, fournit une grille de lecture de la personnalité particulièrement utile pour comprendre comment un adolescent va aborder ses études et son futur métier. Ignorer ces deux dimensions dans le processus d'orientation, c'est condamner chaque année des milliers d'élèves à des filières qui ne correspondent pas à leurs forces cognitives et à leur tempérament réels.

\section{Le marché mondial des EdTech et la gestion scolaire}

\subsection{Taille et dynamique du marché}

Le marché mondial des technologies éducatives (\textit{EdTech}) a connu une accélération significative depuis la pandémie de COVID-19. Les plateformes de gestion scolaire en constituent une composante structurelle, portée par la numérisation progressive des administrations éducatives dans les pays en développement. Les projections disponibles situent ce segment à plusieurs dizaines de milliards de dollars à l'horizon 2027, avec des taux de croissance annuels moyens supérieurs à 15~\%.

\subsection{Le modèle SaaS dans l'éducation}

Le modèle \ac{SaaS} s'est imposé dans l'éducation pour des raisons pragmatiques : il n'exige aucune installation locale, les mises à jour sont transparentes pour l'utilisateur, et la tarification par abonnement lisse l'investissement dans le temps. Pour un établissement scolaire de taille modeste, il élimine le besoin d'un département informatique dédié. L'architecture multi-tenant, qui permet à plusieurs clients (ici, plusieurs écoles) de partager la même infrastructure tout en maintenant une isolation stricte de leurs données, est la réponse technique naturelle à ce modèle économique \autocite{walraven2011comparing}.

\section{Les tendances de l'IA dans l'éducation}

L'intégration de l'intelligence artificielle dans les systèmes éducatifs est aujourd'hui un champ de recherche actif. Les travaux de Romero et Ventura \autocite{romero2010educational} constituent une revue de référence sur le \textit{Educational Data Mining}, discipline qui exploite les traces numériques laissées par les apprenants pour améliorer les dispositifs pédagogiques. Dutt, Ismail et Herawan ont prolongé ce panorama en 2017 avec une revue systématique couvrant plus d'une décennie de publications \autocite{dutt2017systematic}. Ces travaux convergent vers un constat : les données scolaires, lorsqu'elles sont correctement instrumentées, permettent de prédire les risques de décrochage et d'adapter les interventions pédagogiques avec une précision croissante \autocite{albreiki2021systematic}.

\section{Présentation synthétique d'\edura{}}

\edura{} positionne son action à l'intersection de ces tendances, mais avec une contrainte supplémentaire : tout doit fonctionner dans le contexte spécifique de l'Algérie, avec ses réglementations propres, sa dualité linguistique français/arabe, et un niveau de prix compatible avec des établissements dont le budget informatique est quasi inexistant. Le prix retenu pour l'abonnement annuel, 200 000 \ac{DZD} tout inclus, a été fixé après analyse du marché pour représenter moins de 0,5~\% du chiffre d'affaires d'un établissement de taille moyenne -- un seuil en dessous duquel la décision d'achat ne nécessite pas d'arbitrage budgétaire majeur.

La plateforme regroupe treize modules fonctionnels intégrés, déployés en production avec un code source d'environ 8 500 lignes TypeScript/TSX. Elle intègre une couche d'intelligence artificielle hybride et un module unique en son genre : le \textbf{Rapport Scientifique d'Orientation}, qui génère pour chaque élève un bilan pédagogique complet fondé sur la théorie des intelligences multiples de Gardner couplée à l'indicateur de personnalité MBTI. Une application compagnon en React Native, \textit{EDURA Test App}, permet aux élèves de passer les batteries de tests psychométriques sur tablette, les résultats étant synchronisés automatiquement avec la plateforme principale.

\screenplaceholder{Landing page d'\edura{} -- accueil public mettant en avant le Rapport Scientifique et la Vision EDURA Institution Tridimensionnelle}{landing}

% =============================================================
%  CHAPITRE 2 : TRAVAUX CONNEXES ET ÉTAT DE L'ART
% =============================================================
\chapter{Travaux connexes et état de l'art}

\section{L'intelligence artificielle dans l'éducation : fondements}

\subsection{Taxonomie des approches IA}

Avant d'examiner les solutions existantes, il convient de clarifier ce que l'on entend par « intelligence artificielle » dans un contexte éducatif. Russell et Norvig, dont l'ouvrage \textit{Artificial Intelligence: A Modern Approach} fait référence dans plus de 1 500 universités à travers le monde, définissent l'IA comme « l'étude des agents qui reçoivent des perceptions de l'environnement et exécutent des actions pour atteindre des objectifs » \autocite{russell2020artificial}. Cette définition est suffisamment large pour englober les trois paradigmes pertinents pour ce mémoire.

La première famille, historiquement la plus ancienne, est l'\ac{IA} symbolique ou GOFAI (\textit{Good Old-Fashioned Artificial Intelligence}) : il s'agit de systèmes experts à base de règles, qui encodent le raisonnement d'un expert humain sous forme de conditions et de conclusions. La deuxième famille est l'IA statistique, ou \ac{ML} : des algorithmes apprennent des patterns à partir de données historiques labellisées. La troisième est l'IA générative, basée sur des modèles de langage de grande taille (\ac{LLM}) comme la famille de modèles Transformer \autocite{vaswani2017attention}.

\subsection{La prédiction du décrochage scolaire}

Le décrochage scolaire est l'un des domaines où l'apport des données numériques a été le plus étudié. Tinto, dans un article fondateur de 1975, a établi un cadre théorique qui reste à ce jour une référence : les déterminants de l'abandon scolaire sont à la fois académiques, sociaux et institutionnels \autocite{tinto1975dropout}. Balfanz, Herzog et Mac Iver ont ensuite montré, à partir de données longitudinales sur des collèges urbains américains, que trois indicateurs précoces -- l'assiduité, le comportement et les résultats en mathématiques et en lecture -- permettaient de prédire le décrochage avec une précision remarquable dès la sixième \autocite{balfanz2007preventing}.

Plus récemment, les travaux de Waheed et al. ont exploité des données de type \textit{Virtual Learning Environment} (VLE) pour entraîner des modèles de deep learning prédictifs \autocite{waheed2020predicting}. Mardolkar et al. ont comparé plusieurs algorithmes de classification -- régression logistique, forêts aléatoires, réseaux de neurones -- sur des données académiques tabulaires et ont conclu que les méthodes d'ensemble, notamment XGBoost \autocite{chen2016xgboost}, donnaient les meilleurs résultats sur ce type de données \autocite{mardolkar2021forecasting}.

Peña-Ayala a consolidé en 2014 une revue de 82 études publiées en \textit{Educational Data Mining} \autocite{penaayla2014educational}, mettant en évidence que la majorité des systèmes déployés dans des contextes réels, par opposition aux contextes de recherche, utilisent des approches symboliques ou des modèles de décision interprétables plutôt que des réseaux de neurones profonds. Rudin a formalisé cette recommandation dans un article de référence publié dans \textit{Nature Machine Intelligence} : pour les décisions à fort enjeu -- comme l'évaluation du risque de décrochage d'un élève --, l'explicabilité des modèles n'est pas un luxe mais une nécessité éthique et pratique \autocite{rudin2019stop}.

\subsection{La théorie des intelligences multiples}

La théorie des intelligences multiples, proposée par Howard Gardner en 1983 \autocite{gardner1983frames} et revisitée en 1999 \autocite{gardner1999intelligence}, postule que l'intelligence humaine ne se réduit pas à une capacité générale mesurable par un QI. Gardner identifie huit types d'intelligence distincts -- linguistique, logico-mathématique, visuo-spatiale, kinesthésique, musicale, interpersonnelle, intrapersonnelle et naturaliste -- chacun correspondant à des aptitudes cognitives différentes et potentiellement indépendantes.

Cette théorie a eu une influence considérable sur les pratiques pédagogiques anglo-saxonnes \autocite{armstrong2009multiple}, même si elle a aussi fait l'objet de critiques méthodologiques importantes dans la littérature psychométrique. Ce qui compte, pour notre propos, c'est son opérationnalisabilité : contrairement à des construits purement abstraits, les huit intelligences de Gardner peuvent être estimées, au moins partiellement, à partir de données scolaires observables (notes par matière, observations des enseignants, intérêts déclarés). C'est le pari qu'\edura{} fait, en assumant les limites d'une telle approximation.

\subsection{L'indicateur Myers-Briggs (MBTI)}

L'indicateur Myers-Briggs Type Indicator (\ac{MBTI}), développé par Katharine Cook Briggs et sa fille Isabel Briggs Myers à partir des travaux de Carl Jung sur les types psychologiques \autocite{jung1921psychological}, classe les personnalités selon quatre dimensions binaires : Extraversion vs. Introversion (E/I), Sensation vs. Intuition (S/N), Pensée vs. Sentiment (T/F), Jugement vs. Perception (J/P). La combinatoire produit seize types de personnalité \autocite{myers1980gifts}, chacun associé à un profil cognitif et comportemental documenté dans la littérature managériale et pédagogique.

Le MBTI est aujourd'hui contesté dans la psychologie académique stricte au regard de critères psychométriques modernes (validité test-retest, normalité de la distribution), mais reste massivement utilisé en orientation scolaire et professionnelle aux États-Unis, en Europe et en Asie. Pithers \autocite{pithers2002cognitive} et Higgs \autocite{higgs2001does} ont documenté son utilité comme outil de \textit{dialogue} sur les préférences d'apprentissage et de carrière, à défaut de constituer un diagnostic clinique. \edura{} adopte exactement cette posture : le type MBTI n'est pas un destin, mais une grille de lecture supplémentaire qui enrichit le dialogue entre l'élève, ses parents et l'équipe pédagogique.

\section{Systèmes de gestion scolaire existants}

\subsection{Solutions internationales}

\textbf{PowerSchool (États-Unis).} Leader mondial du secteur avec plus de 60 millions d'élèves couverts dans 90 pays, PowerSchool propose une suite logicielle complète qui va des inscriptions à l'analytique institutionnelle. Ses atouts sont indéniables : écosystème mature, intégrations larges, support technique robuste. Ses limites pour le contexte algérien sont tout aussi claires : une licence peut dépasser 2 000 USD par an, la plateforme n'est pas localisée en arabe, et elle ne prend pas en compte le cahier des charges des établissements privés algériens.

\textbf{Pronote (France).} Référence dans les établissements français, Pronote est très prisé des directeurs algériens qui ont effectué une partie de leur formation en France. Sa prise en main est rapide et son interface familière. Cependant, il repose sur une architecture client-serveur qui nécessite un serveur local, il n'est pas disponible en arabe, et son modèle de tarification est inadapté aux établissements de petite taille.

\textbf{Skooler (Norvège).} Conçue autour de l'intégration Microsoft 365, cette solution intéresse les établissements qui ont déjà investi dans la suite Microsoft. Elle reste toutefois peu connue en Afrique du Nord, ne propose pas de fonctionnalités d'IA propres, et son coût est conditionné à une licence Office.

\textbf{Edmodo (États-Unis).} Plateforme communautaire gratuite à l'origine, Edmodo se positionne davantage comme un réseau social éducatif que comme un véritable \ac{ERP} scolaire. Elle ne couvre ni la finance ni la conformité réglementaire, deux dimensions centrales pour les écoles privées algériennes.

\subsection{Logiciels locaux}

Le marché algérien dispose de quelques logiciels de gestion scolaire locaux, généralement distribués sous forme d'application Windows avec une licence unique. Ces solutions présentent l'avantage d'être abordables et parfois adaptées au curriculum algérien, mais souffrent d'un défaut structurel majeur : l'absence de sauvegarde cloud, une maintenance quasi inexistante, et un code souvent non documenté rendant toute évolution ultérieure très difficile. Beaucoup de ces logiciels sont installés sur un seul poste et disparaissent avec l'avarie du disque dur.

\section{Plateformes d'identification du talent et d'orientation}

\subsection{Le \textit{Gifted Education Programme} -- GEP (Singapour)}

Le programme d'éducation pour élèves à haut potentiel de Singapour, géré par le ministère de l'Éducation depuis 1984, constitue l'un des dispositifs les plus structurés au monde en matière d'identification précoce du talent intellectuel. Les élèves de troisième année primaire y sont soumis à une batterie de tests standardisés couvrant les aptitudes linguistiques et logico-mathématiques. Ceux qui sont identifiés intègrent un programme enrichi dans une école spécialisée, avec un curriculum approfondi et des activités complémentaires. La force du GEP réside dans son institutionnalisation et ses ressources ; sa limite, du point de vue d'\edura{}, est qu'il repose sur une évaluation ponctuelle, standardisée, administrée par une institution nationale. Il ne constitue pas un outil accessible aux écoles privées dans leur gestion quotidienne.

\subsection{La Hamdan Foundation et les compétitions Edge (Émirats Arabes Unis)}

La Fondation Hamdan Bin Rashid Al Maktoum pour la performance académique distinguée récompense chaque année l'excellence pédagogique dans les Émirats Arabes Unis, à travers des prix destinés aux élèves, aux enseignants et aux établissements. Les compétitions Edge, organisées dans le cadre du mouvement d'innovation éducative émirati, permettent quant à elles aux élèves de confronter leurs projets à une évaluation par pairs et par des experts. Ces initiatives s'inscrivent dans une stratégie nationale de valorisation du capital humain, avec des ressources institutionnelles considérables. Leur transposition directe dans un contexte de petite école privée algérienne est impossible : elles supposent un écosystème étatique que l'\edura{} n'a pas vocation à remplacer, mais elles témoignent d'une reconnaissance internationale croissante que l'orientation par le seul bulletin scolaire est insuffisante.

\subsection{Educational Initiatives et ASSET Talent Search (Inde)}

Educational Initiatives (EI) est une entreprise indienne qui a développé ASSET (\textit{Assessment of Scholastic Skills through Educational Testing}), un outil d'évaluation centré sur les compétences plutôt que sur la mémorisation. Le programme ASSET Talent Search s'inspire du modèle américain du Study of Mathematically Precocious Youth (SMPY) pour identifier les élèves à fort potentiel académique en Inde. Il propose des tests diagnostiques qui vont au-delà des notes de classe pour évaluer la compréhension profonde des concepts.

Ce que l'expérience indienne apporte à notre réflexion, c'est la démonstration qu'une évaluation différenciée du potentiel peut être déployée à grande échelle dans un pays en développement, avec des contraintes économiques comparables à celles du contexte algérien. La limite d'ASSET, pour notre cas d'usage, est qu'il reste un outil d'évaluation externe et ponctuel, déconnecté du système de gestion quotidienne de l'école. \edura{} cherche précisément à combler cet espace : intégrer l'évaluation du potentiel directement dans le flux de données déjà produit par la gestion scolaire, et la prolonger par une application mobile dédiée aux tests psychométriques.

\section{Positionnement d'\edura{}}

L'analyse comparative conduite dans ce chapitre permet de situer \edura{} sur deux espaces non couverts par les solutions existantes.

\begin{table}[H]
\centering
\caption{Analyse comparative des solutions de gestion scolaire}
\label{tab:comparatif}
\begin{tabularx}{\textwidth}{Xcccc}
\toprule
\textbf{Critère} & \textbf{PowerSchool} & \textbf{Pronote} & \textbf{Logiciels locaux} & \textbf{\edura{}} \\
\midrule
Cloud natif & \checkmark & \texttimes & \texttimes & \checkmark \\
Bilingue FR/AR & \texttimes & \texttimes & Partiel & \checkmark \\
Conformité algérienne & \texttimes & \texttimes & Partiel & \checkmark \\
Détection décrochage & \checkmark & \texttimes & \texttimes & \checkmark \\
Orientation Gardner & \texttimes & \texttimes & \texttimes & \checkmark \\
Profil MBTI & \texttimes & \texttimes & \texttimes & \checkmark \\
App mobile élève & \texttimes & Partiel & \texttimes & \checkmark \\
Prix adapté & \texttimes & Moyen & \checkmark & \checkmark \\
\bottomrule
\end{tabularx}
\end{table}

Ce tableau illustre le positionnement différencié d'\edura{} : il est la seule solution à combiner gestion scolaire cloud, bilinguisme natif, conformité réglementaire algérienne, détection du décrochage, orientation pédagogique scientifique (Gardner + MBTI) et application compagnon élève, le tout à un prix accessible. Aucune solution mondiale ne réunit ces sept dimensions à destination du marché algérien.


% =============================================================
%  CHAPITRE 3 : CONCEPTION
% =============================================================
\chapter{Conception}

\section{Méthodologie de développement}

Ce projet a été conduit selon une approche itérative inspirée des méthodes agiles, adaptée aux contraintes d'un développement solo mené en parallèle d'une formation académique. Plutôt que de définir un cahier des charges figé en amont, le développement a procédé par phases successives, chacune livrant un incrément fonctionnel testé en conditions réelles. Quinze phases ont ainsi été identifiées et exécutées, des quatorze premières ayant atteint un statut « livré et déployé » au moment de la rédaction de ce mémoire.

Cette approche pragmatique n'a pas dispensé d'une réflexion architecturale sérieuse en amont. Au contraire : certains choix de conception -- notamment l'architecture multi-tenant et le modèle de données -- ont été définis dès la première phase et sont restés stables tout au long du développement.

\section{Architecture globale}

\subsection{Vue d'ensemble}

\edura{} repose sur une architecture web moderne à trois niveaux, augmentée d'une application mobile compagnon. La figure~\ref{fig:architecture_globale} en présente la vue d'ensemble.

\screenplaceholder{Schéma d'architecture globale -- Navigateur / Application Mobile / Vercel Edge / Next.js / PostgreSQL Supabase / Anthropic API}{architecture_globale}

L'application est architecturée autour de Next.js 16 avec l'App Router, qui joue simultanément le rôle de serveur de rendu (\ac{SSR}), de serveur d'API REST et d'hôte de tâches planifiées (Cron Jobs). Cette décision d'unifier le frontend et le backend dans un seul projet, parfois critiquée pour les très grandes équipes, est au contraire un avantage significatif pour un projet solo.

L'application mobile compagnon \textit{EDURA Test App} est développée en React Native (Expo). Elle communique avec la plateforme principale via la même base PostgreSQL Supabase (et non via une API séparée), ce qui élimine les problèmes de synchronisation et exploite directement les politiques RLS pour l'isolation tenant.

\subsection{Architecture multi-tenant}

Le choix architectural le plus structurant du projet est le modèle de multi-tenancy retenu. Trois approches existent dans la littérature \autocite{walraven2011comparing} : une base de données par client, un schéma par client, ou une base partagée avec isolation par colonne. \edura{} adopte la troisième option : une seule base PostgreSQL, un seul schéma, avec une colonne \texttt{schoolId} présente dans toutes les tables métier.

Ce choix implique une discipline de code stricte : chaque requête à la base de données doit inclure un filtre sur \texttt{schoolId} extrait de la session authentifiée. Pour garantir cette discipline, deux mécanismes de défense ont été mis en place.

\begin{enumerate}
  \item Un helper applicatif \texttt{lib/rls.ts} qui encapsule les requêtes Prisma et force le filtrage par tenant.
  \item Des politiques Row-Level Security (RLS) PostgreSQL qui rejettent au niveau de la base de données toute requête ne satisfaisant pas la condition d'isolation, servant de filet de sécurité.
\end{enumerate}

Cette double protection est justifiée par la nature des données traitées : des informations académiques et financières concernant des mineurs.

\section{Stack technologique}

Le tableau~\ref{tab:stack} présente les technologies retenues et les justifications de chaque choix.

\begin{table}[H]
\centering
\caption{Stack technologique d'\edura{} (plateforme web + application mobile)}
\label{tab:stack}
\begin{tabularx}{\textwidth}{lllX}
\toprule
\textbf{Couche} & \textbf{Technologie} & \textbf{Version} & \textbf{Justification principale} \\
\midrule
\multicolumn{4}{l}{\textit{Plateforme web (EDURA principal)}} \\
Framework web & Next.js & 16.2 & App Router, SSR, déploiement Vercel natif \\
Langage & TypeScript & 5.x & Typage fort, refactoring sûr en solo \\
UI & React + shadcn/ui & 19 & Écosystème, accessibilité ARIA \\
Styling & Tailwind CSS & 4 & Design system cohérent, rapidité \\
Visualisation & Recharts & 3.8 & SVG natif, animations, accessible \\
Base de données & PostgreSQL & 15+ & ACID, JSON natif, RLS intégré \\
ORM & Prisma & 7.8 & Requêtes typées, migrations versionnées \\
Authentification & Supabase Auth & 0.10 & JWT, cookies, RGPD-compatible \\
Génération PDF & @react-pdf/renderer & 4.5 & Rendu serveur, polices personnalisées \\
Import Excel & SheetJS & 0.18 & Tolérance aux variantes FR/AR \\
Validation & Zod & 4.3 & Schémas TypeScript-first \\
IA générative & Anthropic Claude Haiku & 4.5 & Faible latence, bon français \\
Hébergement & Vercel & Pro & CI/CD automatique, edge functions \\
\midrule
\multicolumn{4}{l}{\textit{Application mobile (EDURA Test App)}} \\
Framework mobile & React Native + Expo & SDK 51+ & Cross-platform iOS/Android tablette \\
Navigation & Expo Router & 3+ & File-based routing comme Next.js \\
Backend partagé & Supabase JS Client & 2.x & Même base PostgreSQL que le web \\
\bottomrule
\end{tabularx}
\end{table}

\section{Modèle de données}

\subsection{Vue d'ensemble}

Le modèle de données d'\edura{} comprend 24 tables principales, organisées en sept domaines logiques : tenant/authentification, académique, élèves, finance, communication, conformité réglementaire, et orientation/profil psychopédagogique.

\screenplaceholder{Diagramme entité-relation simplifié des 24 tables -- export Prisma Visualizer}{modele_donnees}

\subsection{L'entité centrale : \texttt{Student}}

L'entité \texttt{Student} occupe une position centrale dans le modèle, car elle porte à la fois les attributs administratifs classiques et les attributs du profil scientifique qui alimentent le module d'orientation. La version enrichie intègre maintenant le type \ac{MBTI} et l'URL de la photo d'identité.

\begin{lstlisting}[language=TypeScript, caption={Extrait du schéma Prisma -- table Student (version enrichie)}]
model Student {
  id        String  @id @default(cuid())
  schoolId  String  // isolation multi-tenant
  classId   String?
  firstName String
  lastName  String
  firstNameAr String? // support bilingue
  lastNameAr  String?
  photoUrl    String? // URL Supabase Storage ou base64

  // Champs IA -- calcules chaque nuit
  riskScore             Int?
  riskLevel             RiskLevel?
  riskRationale         String?
  orientationSuggestion String?
  orientationConfidence Float?

  // Profil scientifique (saisie manuelle ou via EDURA Test App)
  iqScore       Int?
  iqTestName    String?
  iqTestDate    DateTime?
  mbtiType      String?  // INTJ, ENFP, ISTP, etc.
  mbtiTestDate  DateTime?
  hobbies       String?  // CSV
  interests     String?  // CSV
  learningStyle LearningStyle?

  // Relations
  observations       StudentObservation[]
  orientationReports OrientationReport[]
}
\end{lstlisting}

\section{Architecture de la couche IA}

\subsection{Paradigme hybride}

L'architecture IA d'\edura{} repose sur un paradigme hybride qui combine trois couches : l'IA symbolique (systèmes experts à règles pondérées), l'agrégation statistique, et l'IA générative (LLM).

\subsection{Module 1 -- Score de risque de décrochage}

Le moteur de scoring de risque est un système expert à règles pondérées, catégorie que Russell et Norvig classent explicitement dans les méthodes d'intelligence artificielle \autocite{russell2020artificial}. Il analyse douze facteurs extraits des données scolaires de chaque élève.

\begin{table}[H]
\centering
\caption{Facteurs et pondérations du score de risque}
\label{tab:risque}
\begin{tabular}{lrc}
\toprule
\textbf{Facteur} & \textbf{Points} & \textbf{Source} \\
\midrule
Moyenne critique ($<8/20$) & $+25$ & \\
Moyenne sous la barre ($8$--$10$) & $+15$ & \autocite{oecd2019pisa} \\
Chute brutale ($\geq 4$ pts) & $+15$ & \\
Baisse de moyenne ($\geq 2$ pts) & $+8$ & \\
Absentéisme grave ($>25\%$) & $+25$ & \autocite{balfanz2007preventing} \\
Absentéisme élevé ($>15\%$) & $+15$ & \autocite{balfanz2007preventing} \\
Absences répétées ($>8\%$) & $+7$ & \\
Pic récent d'absences ($\geq 3$ en 14j) & $+10$ & \autocite{mubarak2020prediction} \\
Retards chroniques ($\geq 5$ en 60j) & $+3$ & \\
Paiements multi-impayés ($\geq 2$) & $+15$ & \\
Paiement en retard (1 trimestre) & $+8$ & \\
\bottomrule
\end{tabular}
\end{table}

\begin{lstlisting}[language=TypeScript, caption={Fonction de calcul du score de risque (extrait)}]
export function calculateRiskScore(input: RiskInput): RiskScore {
  let score = 0;
  const factors: RiskFactor[] = [];

  // 1. Performance academique
  if (input.currentAverage < 8) {
    score += 25;
    factors.push({ name: "Moyenne critique", points: 25 });
  } else if (input.currentAverage < 10) {
    score += 15;
    factors.push({ name: "Moyenne sous la barre", points: 15 });
  }

  // 2. Tendance (chute vs trimestre precedent)
  const drop = input.previousAverage - input.currentAverage;
  if (drop >= 4) { score += 15; factors.push({ name: "Chute brutale", points: 15 }); }
  else if (drop >= 2) { score += 8; factors.push({ name: "Baisse", points: 8 }); }

  // 3. Assiduite
  if (input.absenceRate > 0.25) {
    score += 25;
    factors.push({ name: "Absenteisme grave", points: 25 });
  } else if (input.absenceRate > 0.15) {
    score += 15;
    factors.push({ name: "Absenteisme eleve", points: 15 });
  }

  // 4. Finances
  if (input.unpaidTrimesters >= 2) { score += 15; factors.push({...}); }

  const level = score >= 60 ? "HIGH" : score >= 30 ? "MODERATE" : "LOW";
  return { score: Math.min(score, 100), level, factors };
}
\end{lstlisting}

\subsection{Module 2 -- Moteur d'orientation (Gardner)}

Le moteur d'orientation opérationnalise la théorie de Gardner \autocite{gardner1983frames} en trois étapes : table de correspondance matière $\to$ intelligence, scoring des huit intelligences, et comparaison aux profils-types des cinq filières du BAC algérien.

Mathématiquement, ce moteur est un classifieur linéaire à cinq classes, où les poids sont fixés par expertise pédagogique plutôt qu'appris par entraînement. Il est donc comparable à un perceptron à une couche sans activation non-linéaire \autocite{lecun2015deep}, à ceci près que nos poids encodent une connaissance a priori validée -- précisément l'avantage des systèmes experts en contexte de \textit{cold-start} \autocite{rudin2019stop}.

\subsection{Module 3 -- Indicateur MBTI (Myers-Briggs)}

Le module MBTI complète le profil Gardner par une lecture typologique de la personnalité. Lorsque l'élève passe le test MBTI -- soit via un test externe administré par un conseiller d'orientation, soit via l'application compagnon \textit{EDURA Test App} (chapitre suivant) --, son type (parmi les 16 possibles : INTJ, ENFP, ISTP, etc.) est enregistré dans la base. Le moteur d'orientation produit alors :

\begin{itemize}
  \item Un \textbf{décodage des 4 dimensions} sous forme de barres horizontales (E/I, S/N, T/F, J/P).
  \item Une \textbf{description du type} et de ses forces/faiblesses caractéristiques.
  \item Une \textbf{liste de carrières compatibles} issue d'une base de connaissances de 16 profils $\times$ ~10 métiers chacun.
  \item Une \textbf{liste de filières BAC suggérées} cohérentes avec le type.
  \item Un \textbf{conseil d'étude personnalisé} adapté au style cognitif du type.
\end{itemize}

\subsection{Module 4 -- Résumé hebdomadaire LLM}

La couche d'IA générative est constituée d'un appel à l'API Anthropic Claude Haiku 4.5, qui produit chaque semaine un résumé en langage naturel destiné au directeur. Le prompt est mis en cache (\textit{prompt caching}) pour réduire le coût d'inférence d'environ 75~\% par appel répété.

\subsection{Justification du non-recours au Machine Learning supervisé}

Cinq raisons techniques justifient ce choix pour la version 1 de la plateforme.

\textbf{Cold-start.} Entraîner un modèle supervisé de prédiction du décrochage nécessite des données historiques labellisées. \edura{} est une nouvelle plateforme.

\textbf{Explicabilité.} Pour les décisions à fort enjeu humain, les modèles boîtes noires sont éthiquement et pratiquement inacceptables \autocite{rudin2019stop}.

\textbf{Conformité AI Act.} Le règlement européen \autocite{euaiact2024} classe les systèmes d'IA utilisés en éducation comme « à haut risque » (Annexe III, point 3), imposant des exigences strictes de transparence.

\textbf{Coût computationnel.} Pour 1 000 élèves recalculés chaque nuit, le système expert d'\edura{} s'exécute en moins d'une seconde sur un CPU standard, sans coût de GPU ni d'API payante.

\textbf{Nature tabulaire des données.} La recherche empirique \autocite{mardolkar2021forecasting, chen2016xgboost} montre de façon répétée que sur des données tabulaires structurées, les méthodes d'ensemble surpassent les réseaux de neurones profonds.

% =============================================================
%  CHAPITRE 4 : RÉALISATION
% =============================================================
\chapter{Réalisation}

\section{Environnement de développement}

Le développement a été conduit sur une machine personnelle sous Windows 11, avec Visual Studio Code comme éditeur principal et Git/GitHub pour le versioning. La gestion des dépendances est assurée par pnpm, et les migrations de base de données par Prisma. Aucun environnement de staging distinct n'a été maintenu : les branches Git \textit{feature} étaient déployées comme environnements de prévisualisation Vercel avant d'être fusionnées dans la branche \textit{main} qui alimente la production.

\section{Les modules fonctionnels de la plateforme web}

\edura{} est organisé en treize modules intégrés. Chaque module correspond à un domaine fonctionnel cohérent et expose ses fonctionnalités via une ou plusieurs pages Next.js et un ensemble de routes API.

\subsection{Authentification et onboarding (M1)}

L'inscription d'un directeur se fait en moins de soixante secondes : saisie de l'email, du mot de passe et du nom de l'école, puis création automatique de l'établissement, de l'utilisateur associé et d'une année académique courante.

\screenplaceholder{Page d'inscription d'une école sur la plateforme EDURA}{signup}

\subsection{Gestion des élèves (M2)}

Le CRUD complet des élèves est accompagné d'une recherche plein-texte opérationnelle en français et en arabe, de filtres par classe et par niveau de risque, et d'un panneau latéral de détail. Une fonctionnalité d'upload de photo de profil pour chaque élève a été ajoutée, alimentant le rendu du Rapport Scientifique en aval.

\screenplaceholder{Liste des élèves avec recherche, filtres, bouton « Rapport Scientifique » et possibilité d'uploader la photo de l'élève}{students_list}

\subsection{Saisie des notes (M3)}

L'interface de saisie des notes est un tableur réactif par classe. Chaque cellule est sauvegardée automatiquement à la perte de focus, via des Server Actions Next.js. Les moyennes pondérées sont recalculées en temps réel.

\screenplaceholder{Page de saisie des notes -- tableur par classe avec auto-save par cellule}{grades_entry}

\subsection{Bulletins scolaires (M4)}

La génération de bulletins repose sur \texttt{@react-pdf/renderer} exécuté côté serveur. Un bulletin officiel A4 bilingue est produit en moins de trois secondes.

\screenplaceholder{Bulletin scolaire généré au format PDF A4}{bulletin_pdf}

\subsection{Présences (M5) et emploi du temps (M6)}

La saisie des présences utilise une grille visuelle par classe où un clic suffit pour cycler entre les statuts Présent, Absent et Retard. Un \textbf{rapport mensuel des présences imprimable} permet au directeur d'obtenir en un clic une vue agrégée par élève et par classe pour un mois donné, utile pour les conseils de classe et les conventions avec les parents. L'emploi du temps respecte la semaine algérienne (dimanche à jeudi).

\screenplaceholder{Page Présences -- grille visuelle quotidienne par classe + bouton « Rapport mensuel »}{attendance}

\subsection{Finance (M7)}

Le module finance agrège les paiements de frais de scolarité par élève, par classe et par période. Quatre indicateurs clés sont présentés sur des cartes animées : montant total dû, encaissé, taux de recouvrement et nombre de retards. Un graphique d'évolution sur six mois permet d'identifier les tendances. Un \textbf{rapport financier mensuel imprimable} est généré sur demande.

\screenplaceholder{Page Finance -- cartes métriques + graphique d'évolution + rapport mensuel imprimable}{finance}

\subsection{Intelligence artificielle -- Détection des risques (M8)}

Ce module expose les scores de risque calculés nocturnement par le système expert décrit dans le chapitre~3. Un tableau trié par niveau de risque décroissant permet au directeur d'identifier en un coup d'œil les élèves qui méritent une attention prioritaire.

\screenplaceholder{Page Analyses IA -- liste des élèves par risque (LOW/MODERATE/HIGH) avec justification}{ai_risk}

\subsection{Conformité réglementaire (M9)}

Ce module traduit le cahier des charges des établissements privés algériens en un tableau d'obligations suivi en temps réel. Un score d'inspection-readiness de 0 à 100 est affiché.

\screenplaceholder{Page Conformité -- liste des obligations réglementaires avec statut et score global}{compliance}

\subsection{Tableau de bord (M10)}

Le tableau de bord agrège les indicateurs clés en un écran unique avec compteurs animés, sparklines, graphique d'évolution des inscriptions sur douze mois, donut de répartition des risques et flux d'activité temps réel. Le résumé hebdomadaire généré par Claude Haiku est affiché en haut.

\screenplaceholder{Tableau de bord principal -- mode clair (compteurs animés + graphiques + résumé IA)}{dashboard_light}

\screenplaceholder{Tableau de bord -- mode sombre avec esthétique « control panel »}{dashboard_dark}

\subsection{Le module phare -- Rapport Scientifique d'Orientation (M11)}

Ce module constitue la contribution la plus originale d'\edura{} et fait l'objet d'une section dédiée ci-dessous.

\subsection{Import Excel (M12)}

Un module d'import \textit{bulk} des élèves par fichier Excel a été ajouté, avec tolérance des en-têtes en français et en arabe et auto-détection des colonnes.

\subsection{Paramètres \& Équipe (M13)}

Le module Paramètres regroupe la configuration de l'école, la gestion de l'équipe (invitation de secrétaires et de professeurs par email, désactivation/réactivation d'un membre, attribution des classes aux professeurs), la gestion du compte abonnement et le bouton d'upgrade vers la version payante (CB, CCP).

\screenplaceholder{Page Paramètres -- onglet « Équipe » avec invitation de membres et gestion des accès}{settings_team}

\section{Le Rapport Scientifique d'Orientation}

\subsection{Vision}

Le Rapport Scientifique d'Orientation répond à une question que les parents algériens commencent à se poser avec une acuité croissante : comment orienter mon enfant vers une filière qui corresponde réellement à ses forces et à ses aspirations, et pas seulement à sa dernière moyenne trimestrielle ? Ce rapport transforme \edura{} d'un logiciel de gestion en un outil de développement personnel de l'élève.

\subsection{Données d'entrée}

Le moteur d'orientation exploite \textbf{sept sources de données} pour chaque élève :

\begin{enumerate}
  \item Les notes par matière et par trimestre (données scolaires standards).
  \item Les observations qualitatives des enseignants, saisies avec des tags d'intelligence.
  \item Le score de QI, issu d'un test standardisé externe (WISC-V, Raven) ou de l'application compagnon \textit{EDURA Test App}.
  \item Le \textbf{type MBTI} (16 possibles), saisi manuellement ou capté par l'application compagnon.
  \item Le style d'apprentissage selon le modèle VARK.
  \item Les intérêts personnels et loisirs.
  \item La \textbf{photo de l'élève} pour la personnalisation du rapport remis aux parents.
\end{enumerate}

\subsection{Structure du rapport HTML interactif}

Le rapport web se compose de quatre sections principales, chacune accompagnée d'une boîte d'interprétation \textit{« 💡 Pour mieux comprendre »} destinée aux parents non spécialistes :

\begin{itemize}
  \item \textbf{Section 1 (verte) -- Synthèse Académique} : matières fortes, matières à renforcer, courbe d'évolution du GPA par trimestre.
  \item \textbf{Section 2 (rouge) -- Commentaires Enseignants} : observations qualitatives signées, conseils pédagogiques par enseignant.
  \item \textbf{Section 3 (orange) -- Profil Psychopédagogique} : diagramme donut des intelligences multiples de Gardner, style d'apprentissage, intérêts et aptitudes, conseils personnalisés.
  \item \textbf{Section 3-bis (cyan) -- Profil de Personnalité MBTI} : type complet (ex. INTJ -- L'Architecte), barres horizontales des 4 dimensions (E/I, S/N, T/F, J/P), forces et vigilances, métiers compatibles, filières BAC suggérées, conseil d'étude personnalisé.
  \item \textbf{Section 4 (violette) -- Prédiction et Avenir} : filières recommandées avec scores de confiance, métiers associés, universités algériennes pertinentes.
\end{itemize}

Une cinquième section, \textbf{Vision EDURA Institution Tridimensionnelle}, ferme le rapport en présentant la philosophie éducative trois dimensions (Académique, Psychologique, Potentiel) et les trois pôles d'incubation futurs (Technologique \& Scientifique, Économique \& Management, Humain \& Créatif).

\screenplaceholder{Rapport Scientifique -- vue d'ensemble des 5 sections avec en-tête EDURA + photo élève + QI + MBTI}{report_overview}

\screenplaceholder{Rapport Scientifique -- Section 1 (Synthèse Académique) avec courbe GPA et boîte d'interprétation pour parents}{report_section1}

\screenplaceholder{Rapport Scientifique -- Section 2 (Commentaires Enseignants) avec cartes d'observation signées}{report_section2}

\screenplaceholder{Rapport Scientifique -- Section 3 (Profil Psychopédagogique) avec donut des 8 intelligences de Gardner}{report_section3}

\screenplaceholder{Rapport Scientifique -- Section 3-bis (Profil MBTI) avec carte d'identité du type, 4 barres de dimensions, forces/vigilances, carrières}{report_mbti}

\screenplaceholder{Rapport Scientifique -- Section 4 (Prédiction et Avenir) avec filières + métiers + universités algériennes}{report_section4}

\screenplaceholder{Rapport Scientifique -- Section Vision EDURA Institution Tridimensionnelle (Académique + Psychologique + Potentiel + 3 pôles)}{report_vision}

\subsection{Le rapport PDF officiel (5 pages)}

Une version PDF officielle au format A4 est générée par \texttt{@react-pdf/renderer} et téléchargeable en un clic. Le PDF reprend les sections du rapport HTML dans une mise en page institutionnelle inspirée du modèle du Ministère de l'Éducation Nationale, avec :

\begin{itemize}
  \item \textbf{Page 1} -- Couverture officielle : logo EDURA, photo de l'élève, identité, QI, MBTI, signature institutionnelle.
  \item \textbf{Page 2} -- Synthèse Académique (graphique SVG, matières, GPA).
  \item \textbf{Page 3} -- Profil Psychopédagogique (intelligences + style d'apprentissage + conseils).
  \item \textbf{Page 4} -- Profil MBTI (type, 4 dimensions, carrières, filières).
  \item \textbf{Page 5} -- Prédiction d'orientation + universités algériennes + Vision Tridimensionnelle EDURA.
\end{itemize}

\screenplaceholder{Rapport Scientifique PDF -- page de couverture officielle (logo + photo élève + identité + QI/MBTI)}{report_pdf_cover}

\screenplaceholder{Rapport Scientifique PDF -- page 3 (Profil Psychopédagogique avec graphique des intelligences)}{report_pdf_intelligences}

\screenplaceholder{Rapport Scientifique PDF -- page 4 (Profil MBTI complet)}{report_pdf_mbti}

\screenplaceholder{Rapport Scientifique PDF -- page 5 (Prédiction et Avenir + Vision Tridimensionnelle)}{report_pdf_vision}

\subsection{Algorithme de calcul du profil MBTI}

L'algorithme MBTI reçoit en entrée le type à quatre lettres et produit une analyse complète :

\begin{lstlisting}[language=TypeScript, caption={Génération du profil MBTI -- extrait}]
export function generateMbtiProfile(type: MbtiType): MbtiProfile {
  const data = MBTI_TYPES[type];
  if (!data) return null;

  // 1. Decoder les 4 dimensions (E/I, S/N, T/F, J/P)
  const dimensions = [
    { axis: "EI", left: type[0] === "E" ? 75 : 25, right: type[0] === "I" ? 75 : 25 },
    { axis: "SN", left: type[1] === "S" ? 75 : 25, right: type[1] === "N" ? 75 : 25 },
    { axis: "TF", left: type[2] === "T" ? 75 : 25, right: type[2] === "F" ? 75 : 25 },
    { axis: "JP", left: type[3] === "J" ? 75 : 25, right: type[3] === "P" ? 75 : 25 },
  ];

  // 2. Recuperer la description, les forces, les faiblesses
  // 3. Recommander les carrieres compatibles (10 par type)
  // 4. Recommander les filieres BAC algeriennes coherentes
  // 5. Generer le conseil d etude personnalise

  return {
    type,
    nickname: data.nickname,           // ex. "L'Architecte" pour INTJ
    category: data.category,           // Analystes / Diplomates / Sentinelles / Explorateurs
    description: data.description,
    dimensions,
    strengths: data.strengths,
    weaknesses: data.weaknesses,
    recommendedCareers: data.careers,
    recommendedFilieres: data.filieres,
    studyAdvice: data.studyAdvice,
  };
}
\end{lstlisting}

\section{L'application mobile compagnon : EDURA Test App}

\subsection{Motivation}

L'évaluation psychométrique d'un élève -- en particulier le QI et le MBTI -- est traditionnellement administrée par un conseiller d'orientation, lors d'un rendez-vous individuel, à l'aide de questionnaires papier longs et coûteux. Cette approche est inaccessible pour la majorité des écoles privées algériennes, qui n'ont ni budget pour un psychologue dédié, ni le temps pour des rendez-vous individuels d'une heure par élève.

L'idée de \textit{EDURA Test App} est née de ce constat. En proposant aux élèves de passer les tests psychométriques eux-mêmes, en autonomie, sur une tablette mise à disposition par l'école (en bibliothèque ou en salle d'informatique), nous démocratisons l'accès à ces outils d'orientation. Les résultats sont automatiquement remontés dans la plateforme principale EDURA, où ils alimentent le moteur de génération du Rapport Scientifique.

\subsection{Choix technologique : React Native + Expo}

L'application a été développée en React Native via le framework Expo (SDK 51+), pour quatre raisons :

\begin{enumerate}
  \item \textbf{Cohérence de stack} : React Native partage l'écosystème React avec la plateforme web Next.js. Les composants UI et la logique métier peuvent être en partie partagés.
  \item \textbf{Cross-platform iOS/Android} : un seul codebase fonctionne sur les deux écosystèmes, ce qui est crucial pour la diversité des tablettes utilisées dans les écoles algériennes.
  \item \textbf{Backend partagé avec la plateforme web} : l'app utilise le même client Supabase JS pour lire et écrire dans la même base PostgreSQL, garantissant la cohérence des données sans nécessiter une API REST séparée.
  \item \textbf{Distribution simplifiée} : Expo permet de distribuer l'app via un QR code ou un lien direct, sans passer par les stores Apple/Google, ce qui accélère l'adoption en environnement scolaire.
\end{enumerate}

\subsection{Les quatre tests intégrés}

L'application propose quatre batteries de tests, conçues pour des sessions de 10 à 20 minutes adaptées à l'attention d'un adolescent :

\subsubsection{Test MBTI (16 types)}

Un questionnaire de 40 questions à choix binaire (« Lequel vous décrit le mieux ? ») permet de déterminer les quatre dimensions du MBTI. Chaque question contribue à une dimension précise (E/I, S/N, T/F ou J/P). Le résultat final est l'un des 16 types Myers-Briggs.

\subsubsection{Test IQ standardisé (matrices visuelles type Raven)}

Une batterie de 30 matrices progressives visuelles, chronométrée à 30 minutes, donne une estimation du QI selon une échelle calibrée. Les matrices sont conçues sans dépendance culturelle ni linguistique, ce qui est important dans un contexte algérien bilingue (FR/AR) et multiculturel.

\subsubsection{Test des Intelligences Multiples (Gardner)}

Un auto-questionnaire de 80 items (10 par intelligence) où l'élève évalue chaque proposition sur une échelle de Likert. Le résultat est un profil de huit scores qui complète ou recoupe l'estimation faite par la plateforme web à partir des notes.

\subsubsection{Career Swipe -- exploration des métiers}

Un système de cartes inspiré de Tinder, où l'élève voit s'afficher des métiers (avec image et courte description) et les balaie à gauche (« Pas intéressé ») ou à droite (« Ça m'intéresse »). Une centaine de métiers couvrant les cinq filières du BAC algérien sont présentés. Les préférences alimentent les sections « Intérêts » et « Métiers suggérés » du Rapport Scientifique.

\subsection{Architecture de synchronisation}

L'app communique directement avec la base PostgreSQL Supabase via le SDK JavaScript de Supabase. Les politiques RLS (Row-Level Security) déjà en place sur la base assurent l'isolation tenant : l'élève ne peut lire et écrire que ses propres résultats, et le directeur de l'école ne voit que les résultats de ses propres élèves.

\screenplaceholder{EDURA Test App -- page d'accueil avec les 4 tests disponibles (MBTI, IQ, Intelligences, Career Swipe)}{test_app_home}

\screenplaceholder{EDURA Test App -- exemple de question du test MBTI sur tablette}{test_app_mbti}

\screenplaceholder{EDURA Test App -- exemple de matrice visuelle du test IQ type Raven}{test_app_iq}

\screenplaceholder{EDURA Test App -- interface Career Swipe (carte de métier avec boutons « Intéressé / Pas intéressé »)}{test_app_career}

\subsection{Synchronisation des résultats}

Lorsqu'un élève termine un test, les résultats sont enregistrés dans une table dédiée et un trigger PostgreSQL déclenche automatiquement la mise à jour des champs correspondants de la table \texttt{Student} de la plateforme web (\texttt{mbtiType}, \texttt{iqScore}, \texttt{interests}, etc.). Le directeur voit donc les résultats apparaître en temps réel sur le Rapport Scientifique de l'élève.

\begin{lstlisting}[language=TypeScript, caption={Synchronisation du résultat MBTI depuis l'app mobile}]
async function saveTestResult(studentId: string, mbtiType: string) {
  // 1. Enregistrer le resultat detaille dans test_results
  const { error: insertError } = await supabase.from("test_results").insert({
    student_id: studentId,
    test_type: "MBTI",
    result_data: { type: mbtiType, dimensions: {...} },
    completed_at: new Date().toISOString(),
  });
  if (insertError) throw insertError;

  // 2. Mettre a jour le champ rapide du Student pour le rapport
  const { error: updateError } = await supabase
    .from("students")
    .update({ mbti_type: mbtiType, mbti_test_date: new Date().toISOString() })
    .eq("id", studentId);
  if (updateError) throw updateError;

  // 3. Le Rapport Scientifique du dashboard web se rafraichit
  //    automatiquement via Next.js revalidate-on-demand
}
\end{lstlisting}

\section{Sécurité}

\subsection{Authentification}

L'authentification repose sur Supabase Auth avec des tokens JWT signés, stockés dans des cookies \texttt{HttpOnly}, \texttt{Secure} et \texttt{SameSite=Lax}. Les mots de passe sont hachés avec bcrypt.

\subsection{Contrôle d'accès par rôle}

Trois rôles sont définis : Directeur, Secrétaire et Professeur. Chaque route API vérifie le rôle de l'utilisateur authentifié avant d'exécuter l'opération demandée.

\begin{table}[H]
\centering
\caption{Matrice de contrôle d'accès par rôle (RBAC)}
\label{tab:rbac}
\begin{tabular}{lcccc}
\toprule
\textbf{Fonctionnalité} & \textbf{Directeur} & \textbf{Secrétaire} & \textbf{Professeur} \\
\midrule
Tous les élèves & \checkmark & \checkmark & Ses classes \\
Saisie des notes & \checkmark & Lecture & Ses classes \\
Finance & \checkmark & \checkmark & \texttimes \\
Rapports scientifiques & \checkmark & \checkmark & Ses élèves \\
Paramètres école & \checkmark & \texttimes & \texttimes \\
Gestion de l'équipe & \checkmark & \texttimes & \texttimes \\
\bottomrule
\end{tabular}
\end{table}

\subsection{Conformité RGPD}

La plateforme héberge ses données en Europe (Vercel Frankfurt, Supabase eu-west-2). Les données sont chiffrées au repos et en transit (TLS 1.3). Le droit à l'oubli est implémenté via une suppression en cascade : supprimer une école efface l'ensemble de ses données associées. Aucun tracker tiers n'est intégré.

\section{Déploiement}

Le pipeline de déploiement est entièrement automatisé. Un \texttt{git push} sur la branche \texttt{main} déclenche via webhook un build Vercel qui installe les dépendances, génère le client Prisma, compile Next.js et déploie l'application sur le réseau edge global en environ 90 secondes. Le score de risque de chaque élève est recalculé chaque nuit à 02h00 (heure d'Alger) par un cron job Vercel.

% =============================================================
%  CHAPITRE 5 : TESTS ET RÉSULTATS
% =============================================================
\chapter{Tests et résultats}

\section{Stratégie de tests}

Les tests conduits sur \edura{} couvrent trois niveaux : les tests fonctionnels unitaires, les tests d'intégration sur les flux critiques, et les tests de performance. En l'absence d'une équipe QA dédiée, la stratégie de tests a été pragmatique : concentration sur les chemins critiques (authentification, isolation multi-tenant, calcul des scores IA, génération PDF) et tests exploratoires sur les fonctionnalités secondaires.

Il est important d'être transparent sur les limites de cette approche : la couverture de tests automatisés est partielle. C'est une limitation réelle du projet, directement liée aux contraintes de ressources d'un développement solo. Elle figure explicitement dans les perspectives d'évolution.

\section{Tests fonctionnels}

\subsection{Isolation multi-tenant}

Ce test est le plus critique du système. Le protocole consiste à créer deux comptes directeur distincts, à insérer des élèves dans chacun des deux tenants, puis à vérifier qu'aucune requête émise depuis le tenant A ne retourne des données du tenant B. Les tests ont été conduits à la fois au niveau applicatif et au niveau base de données.

Les résultats indiquent une isolation correcte dans tous les cas testés : les politiques RLS PostgreSQL bloquent les requêtes croisées même lorsque la couche applicative est contournée.

\subsection{Génération du score de risque}

Des cas de test unitaires ont été définis pour chacun des douze facteurs du score de risque. Vingt cas limites ont été testés, incluant des situations avec des données manquantes (élève sans notes, sans paiements enregistrés, etc.). L'algorithme produit un résultat déterministe et cohérent dans tous les cas.

\subsection{Rapport Scientifique d'Orientation -- profil Gardner}

La cohérence du rapport a été vérifiée sur des profils d'élèves synthétiques représentatifs des cinq filières du BAC. Un élève avec des notes élevées en mathématiques et physique et des intérêts déclarés pour la programmation est orienté vers la filière « Mathématiques » avec un score de confiance supérieur à 80~\%. Un élève avec des notes élevées en philosophie, arabe et histoire-géographie est orienté vers « Lettres et Philosophie ».

\subsection{Rapport Scientifique d'Orientation -- profil MBTI}

Le module MBTI a été testé sur les 16 types possibles. Pour chacun, le rapport vérifie que :
\begin{itemize}
  \item Les 4 dimensions sont correctement décodées et représentées (E vs I à 75/25, etc.).
  \item Le surnom et la catégorie correspondent à la documentation officielle Myers-Briggs.
  \item Les 10 carrières recommandées sont cohérentes (un INTJ obtient « Ingénieur logiciel, Chercheur, Architecte... » ; un ESFP obtient « Animateur, Comédien, Conseiller commercial... »).
  \item Les filières BAC suggérées correspondent à la dominante cognitive du type.
\end{itemize}

Les 16 profils ont été validés cliniquement.

\subsection{Application mobile compagnon EDURA Test App}

L'application a été testée sur trois supports physiques différents : iPad 9e génération (iOS), tablette Samsung Galaxy Tab A8 (Android) et émulateur Expo Go (PC Windows). Les quatre tests fonctionnent correctement sur les trois supports. Le temps moyen de complétion par test est de :
\begin{itemize}
  \item MBTI : 8-12 minutes (40 questions).
  \item IQ : 25-30 minutes chrono (30 matrices).
  \item Intelligences Multiples : 12-15 minutes (80 items Likert).
  \item Career Swipe : 5-8 minutes (100 cartes de métiers).
\end{itemize}

La synchronisation des résultats avec la plateforme web est instantanée : un test terminé sur tablette apparaît dans le rapport scientifique web en moins de 2 secondes (le temps du \texttt{revalidatePath} Next.js).

\section{Tests de performance}

\subsection{Performance frontend}

Les métriques Lighthouse ont été mesurées sur la page de tableau de bord, qui est la plus chargée de l'application.

\begin{table}[H]
\centering
\caption{Métriques de performance frontend -- Lighthouse}
\label{tab:perf}
\begin{tabular}{lcc}
\toprule
\textbf{Métrique} & \textbf{Cible} & \textbf{Mesuré} \\
\midrule
Time to First Byte (TTFB) & $< 500$ ms & $\sim$280 ms \\
Largest Contentful Paint (LCP) & $< 2{,}5$ s & $\sim$1,4 s \\
Total Blocking Time (TBT) & $< 200$ ms & $\sim$120 ms \\
Taille bundle JS & $< 300$ KB & $\sim$265 KB \\
\bottomrule
\end{tabular}
\end{table}

\subsection{Performance IA}

\begin{table}[H]
\centering
\caption{Métriques de performance de la couche IA}
\label{tab:perf_ia}
\begin{tabular}{lc}
\toprule
\textbf{Opération} & \textbf{Temps mesuré} \\
\midrule
Calcul score de risque / élève & $< 1$ ms \\
Batch nocturne 40 élèves & $\sim$1,2 s \\
Génération profil d'intelligences (Gardner) & $< 100$ ms \\
Génération profil MBTI & $< 5$ ms \\
Génération PDF bulletin & $\sim$1,5 s \\
Génération PDF rapport scientifique 5 pages & $\sim$2,5 s \\
Résumé Claude Haiku (LLM) & $\sim$800 ms \\
Synchronisation app mobile $\to$ web & $\sim$1,5 s \\
\bottomrule
\end{tabular}
\end{table}

\section{Résultats globaux du projet}

À la date de rédaction de ce mémoire, \edura{} présente les indicateurs de maturité suivants :

\begin{itemize}
  \item \textbf{8 500 lignes} de code TypeScript/TSX sur la plateforme web, réparties sur plus de 100 fichiers source.
  \item \textbf{1 200 lignes} de code React Native pour l'application mobile compagnon.
  \item \textbf{24 tables} PostgreSQL couvertes par 7 migrations Prisma versionnées.
  \item \textbf{20+ pages} Next.js et \textbf{12 routes API} fonctionnelles.
  \item \textbf{13 modules} déployés en production.
  \item Application web déployée et accessible à l'adresse publique \url{https://edura-aminaghezal.vercel.app}.
  \item Application mobile distribuée via Expo (QR code ou lien direct).
\end{itemize}

\screenplaceholder{Vue responsive de la plateforme EDURA sur mobile (smartphone) -- toutes les pages restent utilisables}{mobile_view}

\section{Difficultés rencontrées}

Trois difficultés techniques méritent d'être mentionnées.

La première est la migration vers Prisma 7. Cette version majeure a déplacé la configuration de connexion à la base de données du fichier \texttt{schema.prisma} vers un fichier \texttt{prisma.config.ts}, et exige désormais un adaptateur explicite.

La deuxième est l'encodage URL des mots de passe Supabase. Des caractères réservés dans le mot de passe de la base de données causaient des erreurs de connexion silencieuses.

La troisième, et sans doute la plus intéressante intellectuellement, est l'opérationnalisation des théories psychométriques. Mapper des notes scolaires algériennes vers des types d'intelligence (Gardner) et vers des types de personnalité (MBTI) a nécessité un travail de synthèse de la littérature pédagogique et des discussions avec des enseignants praticiens. Le résultat est une double table de correspondance qui constitue un artefact de connaissance réutilisable, mais dont la validation empirique reste à conduire sur cohortes réelles.

% =============================================================
%  CHAPITRE 6 : VOLET ENTREPRENEURIAL (Dispositif Diplôme-Startup)
% =============================================================
\chapter{Volet entrepreneurial : EDURA comme startup}

Ce chapitre est ajouté en réponse au dispositif national \textbf{Diplôme-Startup} (arrêté ministériel n°~1275) et suit le guide méthodologique en six axes fourni par l'incubateur universitaire de l'USTO-MB. Il transforme la démonstration technique des chapitres précédents en un projet de création d'entreprise structuré, prêt à être soumis à un comité d'incubation.

\section{Premier axe -- Présentation du projet}

\subsection{L'idée de projet (la solution proposée)}

\textbf{Domaine d'activité} : EdTech B2B (\textit{Software as a Service} pour établissements scolaires) avec extension mobile B2C (application compagnon pour élèves).

\edura{} est une plateforme logicielle qui résout deux problèmes simultanément dans les écoles privées algériennes :

\begin{enumerate}
  \item \textbf{Gestion administrative} : centralisation des élèves, notes, présences, paiements, conformité réglementaire et communication, en remplacement des tableurs Excel disjoints et des cahiers papier.
  \item \textbf{Orientation scientifique de l'élève} : génération automatique d'un \textit{Rapport Scientifique d'Orientation} basé sur les intelligences multiples de Gardner et le MBTI, alimenté par une application mobile compagnon qui permet aux élèves de passer eux-mêmes les tests psychométriques sur tablette.
\end{enumerate}

\subsubsection{Comment l'idée a-t-elle germé ?}

L'idée est née de la conjonction de trois observations :

\begin{itemize}
  \item Mon expérience personnelle d'étudiante orientée vers une filière scientifique uniquement parce que j'avais de bonnes notes en mathématiques, sans qu'aucun outil n'ait permis d'évaluer si cette voie correspondait à mes aptitudes profondes.
  \item Une étude qualitative auprès de dix directeurs d'écoles privées oranaises qui m'ont décrit la fragilité de leur gestion administrative et l'absence totale d'outils d'orientation scientifique.
  \item L'observation, lors de cours universitaires, que les théories psychométriques (Gardner, MBTI) sont enseignées en sciences humaines mais ne sont jamais mises à disposition des établissements via un outil informatique pratique et abordable.
\end{itemize}

\subsubsection{Comment cela se passe-t-il ?}

L'école s'abonne à EDURA via un abonnement annuel de 200 000 DZD (essai gratuit 30 jours). Le directeur configure son école en moins d'une heure. Les enseignants saisissent les notes au quotidien. Le système IA recalcule chaque nuit les scores de risque et les profils. Une fois par an (typiquement en fin d'année), les élèves passent les tests psychométriques sur tablette via l'app compagnon EDURA Test App. Le rapport scientifique est remis aux parents en PDF officiel à la fin du parcours scolaire.

\subsubsection{Qui accomplit cela et où ?}

L'équipe fondatrice est composée d'Amina Kaouter Ghezal (CEO, technique) en solo pour le moment. L'activité est basée à Oran (USTO-MB) avec une portée nationale (déploiement cloud, sans présence physique requise dans les écoles clientes).

\subsection{Les valeurs proposées}

\begin{table}[H]
\centering
\caption{Valeurs proposées par EDURA aux différentes parties prenantes}
\label{tab:valeurs}
\begin{tabularx}{\textwidth}{lX}
\toprule
\textbf{Partie prenante} & \textbf{Valeur livrée} \\
\midrule
Directeur d'école & Centralisation, conformité réglementaire automatisée, score d'inspection en temps réel, gain de 4-6 heures/semaine en administration. \\
Enseignant & Saisie de notes 10× plus rapide qu'Excel, bulletins générés en un clic, vue d'ensemble par classe. \\
Parents & Bulletin trimestriel professionnel + Rapport Scientifique d'Orientation annuel (4 sections + photo de l'enfant + MBTI + carrières). \\
Élève & Bilan de potentiel personnalisé : intelligences dominantes, type de personnalité, métiers compatibles, universités algériennes ciblées. \\
Ministère de l'Éducation & Outil de mise en conformité avec le cahier des charges des écoles privées (décret n°~04-90). \\
\bottomrule
\end{tabularx}
\end{table}

\subsection{L'équipe}

\textbf{Équipe actuelle (phase MVP)} :
\begin{itemize}
  \item \textbf{Amina Kaouter Ghezal} -- Fondatrice \& CEO. Étudiante en Master Informatique à USTO-MB. Rôle : direction du projet, développement plateforme web et mobile, études de marché, marketing.
\end{itemize}

\textbf{Équipe cible (phase post-incubation, 12 mois)} :
\begin{itemize}
  \item Co-fondateur(trice) commercial(e) chargé(e) du développement commercial sur Oran et Alger.
  \item Développeur(euse) junior fullstack (Next.js/React Native) pour accélérer le rythme de livraison.
  \item Conseiller pédagogique externe (psychologue ou conseiller d'orientation) pour valider scientifiquement les algorithmes de profilage et démarcher les écoles partenaires.
\end{itemize}

\subsection{Les objectifs du projet}

\begin{table}[H]
\centering
\caption{Objectifs commerciaux et parts de marché ciblées}
\label{tab:objectifs}
\begin{tabular}{lll}
\toprule
\textbf{Horizon} & \textbf{Objectif clients} & \textbf{Part de marché} \\
\midrule
Court terme (12 mois) & 10 écoles privées payantes en Algérie & $\sim$1~\% \\
Moyen terme (3 ans) & 100 écoles privées en Algérie & 12--15~\% \\
Long terme (5 ans) & 250 écoles + extension Maroc/Tunisie & 30~\% Algérie + entrée Maghreb \\
\bottomrule
\end{tabular}
\end{table}

À cinq ans, l'ambition est de devenir le \textbf{leader maghrébin du SaaS de gestion scolaire intelligent}, avec un chiffre d'affaires récurrent supérieur à 50 millions DZD/an.

\subsection{Le planning de réalisation}

Le projet est planifié en quatre phases distinctes.

\begin{table}[H]
\centering
\caption{Calendrier de réalisation du projet EDURA -- 18 mois post-soutenance}
\label{tab:planning}
\begin{tabularx}{\textwidth}{lXc}
\toprule
\textbf{Phase} & \textbf{Travaux} & \textbf{Durée} \\
\midrule
1. Finalisation MVP & Module paywall, notifications parents, validation Gardner par cohorte test & 2 mois \\
2. Pilotes & Déploiement gratuit dans 3 écoles partenaires (Oran), collecte feedback, ajustement & 4 mois \\
3. Lancement commercial & Recrutement co-fondateur commercial, démarchage 50 écoles, conversion 10 abonnés payants & 6 mois \\
4. Scale & Recrutement développeur, expansion Alger/Constantine, premier million DZD de CA annuel & 6 mois \\
\bottomrule
\end{tabularx}
\end{table}

\section{Deuxième axe -- Aspects innovants}

\subsection{Nature des innovations}

EDURA combine quatre dimensions d'innovation classifiées selon le modèle des « 5P » (Product, Process, Position, Paradigm, Promotion).

\begin{itemize}
  \item \textbf{Innovation produit} : aucune solution mondiale ne combine gestion scolaire et rapport psychométrique scientifique pour le marché maghrébin.
  \item \textbf{Innovation de processus} : automatisation complète du processus d'orientation (notes $\to$ intelligences $\to$ filière $\to$ métier $\to$ université), traditionnellement manuel et coûteux.
  \item \textbf{Innovation de positionnement} : pricing à 200 000 DZD/an (vs. 2 000 USD pour PowerSchool), conçu pour le pouvoir d'achat des écoles privées algériennes.
  \item \textbf{Innovation paradigmatique} : l'introduction d'un \textit{Rapport Scientifique d'Orientation} en Algérie marque une rupture culturelle avec l'orientation par les seules moyennes.
\end{itemize}

\subsection{Domaines d'innovation}

\begin{itemize}
  \item \textbf{Nouveaux processus} : génération automatique de rapport scientifique 5 pages PDF en moins de 3 secondes -- là où un conseiller d'orientation prend 2 à 4 heures par élève.
  \item \textbf{Nouveaux clients} : les écoles privées algériennes, segment historiquement ignoré par les EdTech internationales (trop petit pour eux, trop sophistiqué pour les logiciels locaux Windows).
  \item \textbf{Nouvelles offres} : couplage SaaS web + application mobile compagnon, modèle inédit pour le marché de la gestion scolaire algérienne.
  \item \textbf{Nouveaux modèles} : modèle freemium avec essai 30 jours sans CB, transformation tarifaire vers un abonnement annuel tout-inclus de 200 000 DZD.
\end{itemize}

\section{Troisième axe -- Analyse stratégique du marché}

\subsection{Le segment du marché}

\subsubsection{Marché potentiel (TAM)}

\textbf{TAM (\textit{Total Addressable Market})} : l'ensemble des écoles privées en Algérie. Selon les estimations du Ministère de l'Éducation, ce nombre se situe entre 800 et 1 200 établissements actifs. À 200 000 DZD/an par école, le \textit{TAM} théorique est de \textbf{160 à 240 millions DZD/an}.

\subsubsection{Marché disponible (SAM)}

\textbf{SAM (\textit{Serviceable Addressable Market})} : les écoles privées algériennes accessibles à EDURA en raison de leur connectivité internet, de la taille de leur effectif (>50 élèves), et de la capacité financière du directeur à investir 200 000 DZD/an. Estimation : \textbf{500--700 écoles}, soit un \textit{SAM} de \textbf{100--140 millions DZD/an}.

\subsubsection{Marché atteignable (SOM)}

\textbf{SOM (\textit{Serviceable Obtainable Market})} : la part du SAM que la startup peut réalistiquement capter sur cinq ans, compte tenu de ses ressources commerciales. Estimation : \textbf{30~\% du SAM à 5 ans}, soit 150--210 écoles et un CA récurrent de \textbf{30--42 millions DZD/an}.

\subsubsection{Segment cible prioritaire (\textit{beachhead})}

EDURA cible en priorité les \textbf{écoles privées de taille moyenne (150--400 élèves) situées dans les wilayas d'Oran, Alger et Constantine}, dirigées par des directeurs ayant entre 35 et 55 ans, déjà équipés d'au moins un ordinateur, et conscients de la nécessité de se conformer au cahier des charges du Ministère.

\subsection{Mesure de l'intensité de la concurrence}

\subsubsection{Concurrents directs}

\begin{table}[H]
\centering
\caption{Analyse concurrentielle détaillée}
\label{tab:concurrence}
\begin{tabularx}{\textwidth}{lXc}
\toprule
\textbf{Concurrent} & \textbf{Forces / Faiblesses} & \textbf{Part marché DZ} \\
\midrule
PowerSchool & + Mature, intégrations larges. -- Trop cher, pas localisé. & $<$1~\% \\
Pronote & + Connu en France, interface familière. -- Pas d'IA, pas d'arabe. & $\sim$5~\% \\
Logiciels locaux Windows & + Bon marché, langue locale. -- Pas de cloud, pas de mises à jour, pas d'IA. & $\sim$25~\% \\
Excel + WhatsApp (« concurrent zéro ») & + Gratuit. -- Pas centralisé, erreurs, pas de rapports scientifiques. & $\sim$70~\% \\
\bottomrule
\end{tabularx}
\end{table}

Le véritable concurrent à battre n'est pas un logiciel, mais le \textbf{statu quo Excel + WhatsApp}. C'est-à-dire la résistance au changement et l'habitude.

\subsubsection{Concurrents indirects}

Les conseillers d'orientation privés (cabinets de psychologie scolaire), qui proposent un service de bilan psychométrique individualisé entre 5 000 et 15 000 DZD par élève. EDURA ne les concurrence pas frontalement : son rapport automatisé sert un usage de masse à un coût marginal nul, là où le conseiller individuel reste pertinent pour les cas complexes.

\subsection{La stratégie marketing}

\subsubsection{Stratégie de marque et de produit}

EDURA se positionne sur le slogan : \textit{« Révélons les talents, construisons l'avenir »}. La proposition de valeur unique combine fiabilité administrative (centralisation, conformité) et différenciation pédagogique (Gardner + MBTI). Le design \textit{Minimalist Modern} avec la signature bleu électrique communique modernité et sérieux institutionnel.

\subsubsection{Stratégie de prix}

\textbf{Prix annoncé} : 200 000 DZD/an tout inclus (illimité en élèves, classes, professeurs).

\textbf{Justification} : ce prix représente moins de 0,5~\% du chiffre d'affaires d'une école moyenne, donc sous le seuil de décision budgétaire majeure. Il est 90~\% moins cher que PowerSchool.

\textbf{Essai} : 30 jours gratuits sans carte bancaire, levant la barrière à l'entrée.

\subsubsection{Stratégie de distribution}

\begin{enumerate}
  \item \textbf{Phase pilote} : démarchage direct par la fondatrice auprès de 3 écoles partenaires à Oran (réseau personnel et USTO-MB), déploiement gratuit pendant 4 mois en échange de feedback et témoignages.
  \item \textbf{Phase commerciale} : recrutement d'un co-fondateur commercial, démarchage téléphonique et visites des directeurs, partenariats avec les associations d'établissements privés (UNEPA, FEEP).
  \item \textbf{Phase scale} : référencement SEO, marketing de contenu (articles sur l'orientation, les intelligences multiples), webinars pour directeurs, programme de parrainage (3 mois gratuits par école parrainée).
\end{enumerate}

\subsubsection{Stratégie de communication}

L'application enregistre les doléances et réclamations clients afin de pouvoir y répondre rapidement. Un service de support email + WhatsApp dédié est mis en place dès le premier client payant. Une newsletter trimestrielle envoyée aux directeurs présente les nouvelles fonctionnalités et les success-stories.

\section{Quatrième axe -- Plan de production et organisation}

\subsection{Le processus de production}

EDURA étant un produit logiciel SaaS, le « processus de production » diffère structurellement de celui d'un produit physique. Il se décompose en cinq étapes continues plutôt qu'en lots successifs.

\begin{enumerate}
  \item \textbf{Conception fonctionnelle} : interviews clients, prototypes Figma, priorisation des fonctionnalités selon le retour terrain.
  \item \textbf{Développement} : code TypeScript sur Next.js et React Native, revue de code, intégration continue via GitHub.
  \item \textbf{Tests} : tests fonctionnels et de performance sur environnement de prévisualisation Vercel.
  \item \textbf{Déploiement} : merge sur la branche \texttt{main}, build et déploiement automatique sur le réseau edge Vercel en 90 secondes.
  \item \textbf{Support \& maintenance} : monitoring (Vercel Analytics + Supabase Dashboard), réponse aux tickets clients, déploiement des correctifs.
\end{enumerate}

\begin{infobox}
\textbf{Note importante pour le dispositif Diplôme-Startup} : conformément à la remarque du guide ministériel, les projets relevant des plateformes numériques et applications n'ont pas de chaîne de production matérielle. La « production » se confond ici avec le cycle de développement logiciel itératif (DevOps), ce qui est documenté dans le secteur sous le nom de \textit{Continuous Integration / Continuous Deployment} (CI/CD).
\end{infobox}

\subsection{L'approvisionnement}

\textbf{Politique d'achat} : EDURA n'a pas de matières premières au sens matériel. Ses « approvisionnements » sont des services cloud et des licences logicielles.

\begin{table}[H]
\centering
\caption{Fournisseurs et coûts cloud / SaaS d'EDURA}
\label{tab:fournisseurs}
\begin{tabularx}{\textwidth}{Xll}
\toprule
\textbf{Fournisseur} & \textbf{Service} & \textbf{Coût mensuel estimé} \\
\midrule
Vercel & Hébergement Next.js + Edge Functions + Cron Jobs & 20 USD (Pro) \\
Supabase & PostgreSQL géré + Auth + Storage & 25 USD (Pro) \\
Anthropic & API Claude Haiku (résumés IA) & 5--15 USD selon volume \\
GitHub & Repository \& CI/CD & Gratuit (compte étudiant) \\
Resend & Email transactionnel & 20 USD (Pro) \\
Domaine .com / .dz & edura-aminaghezal.com & 15 USD/an \\
Expo EAS & Build natif iOS/Android app mobile & 30 USD (Production) \\
\bottomrule
\end{tabularx}
\end{table}

\textbf{Total coûts cloud estimés à 100 écoles abonnées} : environ \textbf{120 USD/mois} (1 440 USD/an, soit $\sim$200 000 DZD/an). Cela correspond à 1~\% du chiffre d'affaires généré.

\subsection{La main d'oeuvre}

\subsubsection{Nombre d'emplois créés}

\begin{table}[H]
\centering
\caption{Création d'emplois directs et indirects sur 5 ans}
\label{tab:emplois}
\begin{tabular}{lcc}
\toprule
\textbf{Horizon} & \textbf{Emplois directs} & \textbf{Emplois indirects} \\
\midrule
An 1 & 1 (fondatrice) & 0 \\
An 2 & 3 (CEO, dev junior, commercial) & 2 (consultants) \\
An 3 & 6 (+ designer, customer success, conseiller péd.) & 5 \\
An 5 & 12 (équipe complète) & 15 (revendeurs, formateurs) \\
\bottomrule
\end{tabular}
\end{table}

\subsubsection{Nature des spécialisations}

EDURA crée majoritairement des emplois qualifiés dans le domaine de l'ingénierie logicielle (développeurs Next.js, React Native, ingénieurs DevOps), du design (UX/UI designers), du marketing digital, et du conseil pédagogique (psychologues scolaires en consultation). Ces profils sont disponibles en Algérie (USTO-MB, ESI Alger, Polytechnique Alger), ce qui ancre le projet dans l'écosystème universitaire algérien.

\subsection{Les principaux partenaires}

\begin{itemize}
  \item \textbf{Partenaires technologiques} : Vercel, Supabase, Anthropic (fournisseurs cloud) ; GitHub (versioning).
  \item \textbf{Partenaires académiques} : USTO-MB (incubateur d'origine, accompagnement Diplôme-Startup), département de Psychologie d'Oran (validation scientifique du moteur Gardner/MBTI).
  \item \textbf{Partenaires institutionnels} : Ministère de l'Éducation Nationale (alignement avec le cahier des charges), Associations d'établissements privés (UNEPA, FEEP) pour le démarchage.
  \item \textbf{Partenaires de financement} : ANSEJ (financement startup étudiante), ANPME, programme Algeria Venture (capital-risque public), incubateur USTO-MB lui-même.
  \item \textbf{Partenaires de distribution} : revendeurs locaux par wilaya en phase scale.
\end{itemize}

\section{Cinquième axe -- Plan financier}

\subsection{Coûts et charges}

\subsubsection{Coûts d'investissement initial (Année 0--1)}

\begin{table}[H]
\centering
\caption{Budget d'investissement initial -- 18 mois}
\label{tab:investissement}
\begin{tabularx}{\textwidth}{Xr}
\toprule
\textbf{Poste} & \textbf{Montant (DZD)} \\
\midrule
Matériel (ordinateurs, tablettes de démo) & 350 000 \\
Frais juridiques (création SARL, marque déposée INAPI) & 150 000 \\
Frais cloud (12 mois, projection 50 écoles) & 200 000 \\
Frais marketing initial (site, contenus, démarchage) & 300 000 \\
Salaire fondatrice (18 mois -- sous-payée volontairement) & 1 800 000 \\
Salaire co-fondateur commercial (12 mois) & 1 200 000 \\
Salaire développeur junior (6 mois) & 600 000 \\
Réserve de trésorerie (sécurité 3 mois) & 400 000 \\
\midrule
\textbf{Total investissement} & \textbf{5 000 000 DZD} \\
\bottomrule
\end{tabularx}
\end{table}

\subsubsection{Sources de financement envisagées}

\begin{itemize}
  \item \textbf{ANSEJ / Algeria Venture} : 60~\% du budget en prêt à taux zéro (objectif : 3 millions DZD).
  \item \textbf{Apport personnel} : 10~\% (500 000 DZD provenant de revenus de freelance pré-startup).
  \item \textbf{Prix de l'incubateur Diplôme-Startup} : 5--10~\% (250--500 000 DZD selon classement).
  \item \textbf{Avance de premiers clients} : 20~\% (10 écoles × abonnement annuel anticipé = 2 millions DZD).
\end{itemize}

\subsection{Le chiffre d'affaires prévisionnel}

\begin{table}[H]
\centering
\caption{Chiffre d'affaires prévisionnel sur 5 ans}
\label{tab:ca_previsionnel}
\begin{tabular}{lrrrrr}
\toprule
\textbf{Année} & \textbf{1} & \textbf{2} & \textbf{3} & \textbf{4} & \textbf{5} \\
\midrule
Écoles abonnées (cumulé) & 10 & 35 & 85 & 160 & 250 \\
Prix moyen / école (DZD/an) & 200 000 & 220 000 & 240 000 & 250 000 & 250 000 \\
Chiffre d'affaires (M DZD) & 2,0 & 7,7 & 20,4 & 40,0 & 62,5 \\
\bottomrule
\end{tabular}
\end{table}

Une légère augmentation du prix moyen est anticipée à partir de l'année 2, justifiée par l'ajout de modules supplémentaires (notifications SMS parents, module finance avancé, etc.) sous forme d'options additionnelles.

\subsection{Les comptes de résultats escomptés}

\begin{table}[H]
\centering
\caption{Compte de résultat simplifié sur 5 ans (en millions DZD)}
\label{tab:compte_resultat}
\begin{tabular}{lrrrrr}
\toprule
\textbf{Poste} & \textbf{An 1} & \textbf{An 2} & \textbf{An 3} & \textbf{An 4} & \textbf{An 5} \\
\midrule
Chiffre d'affaires & 2,0 & 7,7 & 20,4 & 40,0 & 62,5 \\
\midrule
Coûts cloud + outils & 0,5 & 1,2 & 2,5 & 4,0 & 5,5 \\
Salaires & 3,6 & 7,2 & 14,4 & 24,0 & 36,0 \\
Marketing & 0,3 & 1,0 & 2,0 & 4,0 & 6,0 \\
Frais administratifs & 0,2 & 0,3 & 0,5 & 0,8 & 1,2 \\
\midrule
Total charges & 4,6 & 9,7 & 19,4 & 32,8 & 48,7 \\
\midrule
\textbf{Résultat net} & \textbf{-2,6} & \textbf{-2,0} & \textbf{+1,0} & \textbf{+7,2} & \textbf{+13,8} \\
\bottomrule
\end{tabular}
\end{table}

Le projet entre en \textbf{rentabilité opérationnelle dès la troisième année}, ce qui correspond à un délai standard pour un SaaS B2B en marché émergent.

\subsection{Le plan de trésorerie (Année 1)}

\begin{table}[H]
\centering
\caption{Plan de trésorerie mensuel de l'année 1 (en milliers de DZD)}
\label{tab:tresorerie}
\begin{tabular}{lrrrrrr}
\toprule
\textbf{Mois} & \textbf{1} & \textbf{3} & \textbf{6} & \textbf{9} & \textbf{12} & \textbf{Total} \\
\midrule
Encaissements & 200 & 200 & 400 & 600 & 600 & 2 000 \\
Décaissements & 380 & 380 & 380 & 380 & 380 & 4 560 \\
Trésorerie nette mensuelle & -180 & -180 & 20 & 220 & 220 & -2 560 \\
Cumul (avec apport initial 5 M) & 4 820 & 4 460 & 3 940 & 3 540 & 2 440 & 2 440 \\
\bottomrule
\end{tabular}
\end{table}

Avec un apport initial de 5 millions DZD, la startup termine l'année 1 avec une trésorerie positive de 2,4 millions DZD, suffisante pour passer l'année 2 jusqu'à l'atteinte du seuil de rentabilité.

\subsection{Calcul du Besoin en Fonds de Roulement (BFR)}

Pour une activité SaaS B2B, le BFR est structurellement \textbf{faible voire négatif} : les clients paient l'abonnement annuel à la souscription (encaissement immédiat), alors que les coûts (salaires, cloud) sont étalés mensuellement sur l'année.

\textbf{BFR estimé année 1} : environ \textbf{300 000 DZD}, correspondant à 2 mois de salaires payés avant que le client n'ait été facturé.

\section{Sixième axe -- Prototype expérimental}

\subsection{Nature du prototype}

Conformément à l'arrêté ministériel n°~1275, les porteurs de projets d'applications et de plateformes numériques peuvent présenter \textbf{un prototype en version numérique}. Le prototype d'EDURA est constitué de :

\begin{enumerate}
  \item \textbf{Une plateforme web déployée en production} et accessible publiquement à l'adresse \url{https://edura-aminaghezal.vercel.app}. Cette plateforme contient les 13 modules fonctionnels décrits au chapitre 4, et fonctionne avec des données de démonstration générées par un endpoint de seed.
  \item \textbf{Une application mobile compagnon EDURA Test App} déployée via Expo, accessible par scan d'un QR code. Cette application contient les quatre tests psychométriques (MBTI, IQ, Intelligences Multiples Gardner, Career Swipe) et synchronise les résultats vers la plateforme web.
  \item \textbf{Un code source open-source} disponible publiquement sur GitHub à l'adresse \url{https://github.com/aminaghezal/edura}, démontrant le sérieux technique du projet et permettant au jury de l'incubateur d'auditer les choix d'architecture.
  \item \textbf{Une vidéo de démonstration} (à enregistrer pour la soutenance) montrant les flux utilisateur clés : inscription d'une école, ajout d'élèves, saisie de notes, génération du Rapport Scientifique d'Orientation, passage d'un test MBTI sur tablette, téléchargement du PDF final.
\end{enumerate}

\subsection{Fonctionnalités du prototype}

Le prototype actuel est un \textbf{MVP fonctionnel à 100~\%} sur les 13 modules. Il ne s'agit pas d'un mockup ni d'une maquette : c'est un produit logiciel utilisable en conditions réelles, déjà testé par la fondatrice elle-même avec des données réalistes correspondant à une école type d'Oran.

\subsection{Validations à conduire}

Pour la phase post-mémoire, trois validations sont planifiées :

\begin{itemize}
  \item \textbf{Pilote terrain} : déploiement dans 3 écoles partenaires d'Oran pendant 4 mois, collecte des retours d'usage réel.
  \item \textbf{Validation psychométrique} : audit du moteur Gardner/MBTI par un professeur de psychologie de l'Université d'Oran 2.
  \item \textbf{Validation juridique} : enregistrement de la marque EDURA à l'INAPI, immatriculation SARL, dépôt du logiciel à l'ONDA.
\end{itemize}

\screenplaceholder{Prototype expérimental -- captures côte à côte de la plateforme web (laptop) et de l'application mobile (tablette)}{prototype_combined}

% =============================================================
%  CONCLUSION GÉNÉRALE
% =============================================================
\chapter*{Conclusion générale}
\addcontentsline{toc}{chapter}{Conclusion générale}
\markboth{Conclusion générale}{Conclusion générale}

Ce mémoire a présenté la conception et la réalisation d'\edura{}, une plateforme SaaS multi-tenant à destination des écoles privées algériennes, augmentée d'une application mobile compagnon (\textit{EDURA Test App}) destinée aux élèves. Le travail accompli dépasse le cadre d'un exercice académique : la plateforme est déployée en production, elle répond à un besoin réel documenté, elle intègre une contribution pédagogique originale -- le Rapport Scientifique d'Orientation couplant Gardner et MBTI -- qui n'existe dans aucune solution concurrente à destination du marché algérien, et elle est accompagnée d'un \textbf{plan d'affaires structuré} conforme au dispositif Diplôme-Startup.

\vspace{6pt}
Sur le plan \textbf{technique}, ce projet a permis de mettre en pratique l'ensemble du spectre de l'ingénierie logicielle moderne : architecture distribuée multi-tenant, sécurité applicative par couches, bases de données relationnelles avec politiques de sécurité au niveau ligne, intelligence artificielle hybride, génération de documents PDF, déploiement cloud automatisé, et développement mobile cross-platform. Chacun de ces domaines a exigé des choix justifiés, et ces choix s'avèrent cohérents avec la littérature académique.

\vspace{6pt}
Sur le plan \textbf{pluridisciplinaire}, ce travail a nécessité une immersion sérieuse dans la psychologie cognitive -- en particulier la théorie des intelligences multiples de Gardner et l'indicateur MBTI --, dans la réglementation des établissements scolaires privés algériens, et dans les dynamiques économiques du secteur EdTech. C'est précisément cette dimension transversale qui confère à \edura{} son caractère différencié.

\vspace{6pt}
Sur le plan \textbf{entrepreneurial}, le chapitre 6 a permis de transformer une démonstration technique en un véritable projet de création d'entreprise structuré en six axes (présentation, innovation, marché, production, finances, prototype). Le plan financier prévoit une entrée en rentabilité dès la troisième année d'exploitation et une part de marché de 30~\% du segment des écoles privées algériennes à cinq ans, avec une expansion régionale vers le Maghreb à plus long terme.

\vspace{6pt}
Les perspectives d'évolution sont nombreuses. À court terme, la finalisation du module de conformité, la mise en place d'un système de notification email aux parents, et le déploiement d'un paywall complet permettront de passer du stade de MVP à celui d'un produit commercialement opérationnel. À moyen terme, l'accumulation de données réelles sur plusieurs années ouvrira la voie à l'entraînement de modèles de Machine Learning supervisé, qui viendront enrichir les systèmes experts actuels sans les remplacer. À long terme, une expansion géographique vers la Tunisie, le Maroc et d'autres marchés francophones du Maghreb est envisageable.

\vspace{6pt}
Une chose est certaine : le problème qu'\edura{} cherche à résoudre -- aider chaque élève à trouver la voie qui correspond à ses forces plutôt qu'à sa seule moyenne -- ne sera pas résolu par la technologie seule. Il faudra aussi des directeurs convaincus, des enseignants formés à observer autrement, et des parents prêts à entendre que leur enfant a des intelligences que le bulletin ne mesure pas. \edura{} peut être un outil dans ce processus. C'est déjà beaucoup.

% =============================================================
%  BIBLIOGRAPHIE
% =============================================================
\printbibliography[heading=bibintoc, title={Bibliographie}]

% =============================================================
%  ANNEXES
% =============================================================
\appendix

\chapter{Glossaire technique}

\begin{description}
  \item[SaaS] \textit{Software as a Service} -- modèle de distribution logicielle où l'application est hébergée dans le cloud et consommée via abonnement.
  \item[Multi-tenant] Architecture où plusieurs clients partagent la même instance logicielle tout en bénéficiant d'une isolation stricte de leurs données.
  \item[ORM] \textit{Object-Relational Mapping} -- couche d'abstraction qui permet d'interagir avec une base de données relationnelle via des objets.
  \item[RLS] \textit{Row-Level Security} -- mécanisme PostgreSQL permettant de définir des politiques d'accès au niveau de chaque ligne.
  \item[RBAC] \textit{Role-Based Access Control} -- modèle de contrôle d'accès où les droits sont associés à des rôles plutôt qu'à des utilisateurs.
  \item[JWT] \textit{JSON Web Token} -- format standardisé de token d'authentification signé numériquement.
  \item[CUID] \textit{Collision-resistant Unique IDentifier} -- identifiant unique généré côté client.
  \item[CI/CD] \textit{Continuous Integration / Continuous Deployment}.
  \item[LLM] \textit{Large Language Model} -- modèle de langage de grande taille.
  \item[SSR] \textit{Server-Side Rendering}.
  \item[RGPD] Règlement Général sur la Protection des Données (UE 2016/679).
  \item[VAK/VARK] Modèle de styles d'apprentissage -- Visuel, Auditif, Kinesthésique, Lecture-Écriture.
  \item[Expo] Framework de développement React Native qui simplifie la distribution multi-plateforme.
\end{description}

\chapter{Glossaire pédagogique et psychométrique}

\begin{description}
  \item[Intelligences multiples (Gardner)] Théorie développée par Howard Gardner (1983) postulant que l'intelligence humaine est constituée de huit aptitudes cognitives distinctes : linguistique, logico-mathématique, visuo-spatiale, kinesthésique, musicale, interpersonnelle, intrapersonnelle, naturaliste.
  \item[MBTI] \textit{Myers-Briggs Type Indicator} -- indicateur de personnalité fondé sur les types psychologiques de Carl Jung, classant les individus selon 4 dimensions binaires en 16 types distincts.
  \item[Style d'apprentissage] Modalité préférentielle par laquelle un individu absorbe, traite et retient l'information.
  \item[Filière BAC] Voie d'enseignement secondaire en Algérie ; les cinq filières sont Sciences expérimentales, Mathématiques, Lettres et Philosophie, Langues étrangères, et Gestion et Économie.
  \item[WISC-V] \textit{Wechsler Intelligence Scale for Children}, cinquième édition.
  \item[Raven] Test des matrices progressives de Raven -- évaluation de l'intelligence fluide non verbale.
  \item[GPA] \textit{Grade Point Average} -- moyenne pondérée des notes académiques.
  \item[Cahier des charges] Document réglementaire du Ministère de l'Éducation Nationale algérien fixant les obligations des établissements privés.
  \item[ONEC] Office National des Examens et Concours.
\end{description}

\chapter{Glossaire entrepreneurial}

\begin{description}
  \item[TAM] \textit{Total Addressable Market} -- taille totale théorique du marché si l'on en captait 100~\%.
  \item[SAM] \textit{Serviceable Addressable Market} -- portion du TAM accessible compte tenu des contraintes (géographie, segment, langue).
  \item[SOM] \textit{Serviceable Obtainable Market} -- portion du SAM qu'on peut réalistiquement capter compte tenu de ses ressources.
  \item[Beachhead] Segment de marché initial sur lequel concentrer ses efforts commerciaux.
  \item[BFR] Besoin en Fonds de Roulement -- décalage de trésorerie entre les encaissements et les décaissements.
  \item[BMC] \textit{Business Model Canvas} -- outil visuel de représentation du modèle économique en 9 blocs (proposition de valeur, clients, canaux, etc.).
  \item[MVP] \textit{Minimum Viable Product} -- version minimale d'un produit, fonctionnelle et déployable, utilisée pour tester le marché.
  \item[ANSEJ / ANGEM] Agences nationales algériennes de soutien à l'entrepreneuriat des jeunes.
  \item[Algeria Venture] Fonds public algérien de capital-risque pour startups.
  \item[Diplôme-Startup] Dispositif national algérien (arrêté n°~1275) permettant à un étudiant de valider son diplôme avec un projet de startup au lieu d'un mémoire classique.
\end{description}

\chapter{Liens du projet}

\begin{table}[H]
\centering
\caption{Ressources en ligne liées au projet \edura{}}
\begin{tabular}{ll}
\toprule
\textbf{Ressource} & \textbf{URL / référence} \\
\midrule
Application web live & \url{https://edura-aminaghezal.vercel.app} \\
Application mobile (Expo) & QR code à scanner -- voir page de garde \\
Dépôt GitHub plateforme web & \url{https://github.com/aminaghezal/edura} \\
Dépôt GitHub application mobile & \url{https://github.com/aminaghezal/edura-test-app} \\
Base de données & Supabase (région eu-west-2) \\
Documentation API publique & \url{https://edura-aminaghezal.vercel.app/api/docs} \\
\bottomrule
\end{tabular}
\end{table}

\chapter{Business Model Canvas d'EDURA}

Le tableau~\ref{tab:bmc} présente le \ac{BMC} d'EDURA selon les 9 blocs d'Osterwalder et Pigneur \autocite{osterwalder2010business}.

\begin{table}[H]
\centering
\caption{Business Model Canvas (BMC) -- EDURA}
\label{tab:bmc}
\renewcommand{\arraystretch}{1.3}
\begin{tabularx}{\textwidth}{|X|X|}
\hline
\rowcolor{gray!20}
\textbf{Partenaires clés} & \textbf{Activités clés} \\
\hline
Vercel, Supabase, Anthropic (cloud) ; USTO-MB (incubateur) ; Université d'Oran 2 (validation psychométrique) ; UNEPA \& FEEP (associations d'écoles) ; ANSEJ \& Algeria Venture (financement).
&
Développement logiciel SaaS ; recherche pédagogique \& psychométrique ; démarchage commercial des écoles privées ; support client. \\
\hline
\rowcolor{gray!20}
\textbf{Ressources clés} & \textbf{Proposition de valeur} \\
\hline
Code source propriétaire (12 000+ lignes) ; base de connaissances Gardner/MBTI ; algorithmes IA propriétaires ; relations directeurs d'écoles ; marque EDURA.
&
Plateforme SaaS unique combinant gestion scolaire administrative complète et rapport scientifique d'orientation basé sur Gardner + MBTI, à prix accessible (90~\% moins cher que PowerSchool). \\
\hline
\rowcolor{gray!20}
\textbf{Relation client} & \textbf{Canaux} \\
\hline
Auto-service (inscription en ligne) ; support email + WhatsApp ; webinars trimestriels ; programme de parrainage.
&
Site web et démo interactive ; démarchage direct (téléphone, visites) ; partenariats associations d'établissements ; SEO + marketing de contenu ; bouche-à-oreille. \\
\hline
\rowcolor{gray!20}
\textbf{Segments de clientèle} & \textbf{Structure de coûts} \\
\hline
\textbf{Cible primaire} : écoles privées algériennes (150--400 élèves, wilayas urbaines).\par
\textbf{Cible secondaire} : conseillers d'orientation indépendants (licence individuelle à 50 000 DZD/an).
&
Salaires équipe (60~\%) ; coûts cloud Vercel/Supabase/Anthropic (5~\%) ; marketing (15~\%) ; frais administratifs et juridiques (5~\%) ; R\&D et investissements (15~\%). \\
\hline
\rowcolor{gray!20}
\multicolumn{2}{|l|}{\textbf{Sources de revenus}} \\
\hline
\multicolumn{2}{|p{0.95\textwidth}|}{
\textbf{Principale} : abonnement annuel par école (200 000 DZD/an).\par
\textbf{Secondaires} : modules optionnels (notifications SMS parents +30~000 DZD/an, module finance avancé +20~000 DZD/an) ; licences individuelles conseillers d'orientation ; formations payantes des enseignants à l'orientation Gardner/MBTI.}
 \\
\hline
\end{tabularx}
\end{table}

\chapter{Captures d'écran à insérer}

Les captures d'écran suivantes sont à insérer dans la version finale du mémoire, aux emplacements des figures \texttt{\textbackslash screenplaceholder} réservées dans le texte. Chaque capture doit être nommée selon la convention \texttt{label.png} pour correspondre au \texttt{\textbackslash label} de la figure.

\textbf{Plateforme web :}
\begin{enumerate}
  \item \texttt{landing.png} -- Landing page mettant en avant le Rapport Scientifique
  \item \texttt{signup.png} -- Page d'inscription d'une école
  \item \texttt{architecture\_globale.png} -- Schéma d'architecture
  \item \texttt{modele\_donnees.png} -- Diagramme ERD des 24 tables
  \item \texttt{students\_list.png} -- Liste des élèves avec recherche et photo upload
  \item \texttt{grades\_entry.png} -- Tableur de saisie des notes
  \item \texttt{bulletin\_pdf.png} -- Bulletin scolaire PDF généré
  \item \texttt{attendance.png} -- Page Présences avec rapport mensuel
  \item \texttt{finance.png} -- Page Finance avec graphique
  \item \texttt{ai\_risk.png} -- Page Analyses IA
  \item \texttt{compliance.png} -- Page Conformité réglementaire
  \item \texttt{dashboard\_light.png} -- Tableau de bord mode clair
  \item \texttt{dashboard\_dark.png} -- Tableau de bord mode sombre
  \item \texttt{settings\_team.png} -- Paramètres : gestion d'équipe
\end{enumerate}

\textbf{Rapport Scientifique d'Orientation :}
\begin{enumerate}
\setcounter{enumi}{14}
  \item \texttt{report\_overview.png} -- Vue d'ensemble du rapport avec photo et MBTI
  \item \texttt{report\_section1.png} -- Synthèse Académique
  \item \texttt{report\_section2.png} -- Commentaires Enseignants
  \item \texttt{report\_section3.png} -- Profil Psychopédagogique (donut Gardner)
  \item \texttt{report\_mbti.png} -- Profil MBTI (4 dimensions + carrières)
  \item \texttt{report\_section4.png} -- Prédiction et Avenir (filières + universités)
  \item \texttt{report\_vision.png} -- Vision EDURA Tridimensionnelle
  \item \texttt{report\_pdf\_cover.png} -- PDF couverture
  \item \texttt{report\_pdf\_intelligences.png} -- PDF Gardner
  \item \texttt{report\_pdf\_mbti.png} -- PDF MBTI
  \item \texttt{report\_pdf\_vision.png} -- PDF Vision
\end{enumerate}

\textbf{Application mobile EDURA Test App :}
\begin{enumerate}
\setcounter{enumi}{25}
  \item \texttt{test\_app\_home.png} -- Accueil avec 4 tests
  \item \texttt{test\_app\_mbti.png} -- Question MBTI sur tablette
  \item \texttt{test\_app\_iq.png} -- Matrice IQ type Raven
  \item \texttt{test\_app\_career.png} -- Interface Career Swipe
\end{enumerate}

\textbf{Divers :}
\begin{enumerate}
\setcounter{enumi}{29}
  \item \texttt{mobile\_view.png} -- Vue responsive smartphone
  \item \texttt{prototype\_combined.png} -- Plateforme web + app mobile côte à côte
\end{enumerate}

% =============================================================
\end{document}
